\chapter{Mechanizing Floating-Point Bitblasting}
\label{ch:floating-point}
\epigraph{We can be grateful that so useful a work of reference has been created from a project of such awe-inspiring futility}{Robert Poole, Review of Calendrical Calculations}

%%%%%%%%%%%%%%%%%%%%%%%%%%%%%%%%%%%%%%%%%%%%%%%%%%%%%%%%%%%%%%%%%%%%%%%%%%%%%%%
%% Chapter-local macros, ported from the POPL'27 paper preamble.
%%%%%%%%%%%%%%%%%%%%%%%%%%%%%%%%%%%%%%%%%%%%%%%%%%%%%%%%%%%%%%%%%%%%%%%%%%%%%%%

% Number-notation macros: NaN, base-10 (...)_10, base-2 (...)_2,
% bit-pattern <...>_2, and the pack/unpack/classification operators.
\newcommand{\nan}{\ensuremath{\mathrm{NaN}}}
\newcommand{\baseten}[1]{\ensuremath{(#1)_{10}}}
\newcommand{\basetwo}[1]{\ensuremath{(#1)_{2}}}
\newcommand{\bitpattern}[1]{\ensuremath{\langle #1 \rangle_{2}}}
\newcommand{\pack}{\ensuremath{\mathrm{pack}}}
\newcommand{\unpack}{\ensuremath{\mathrm{unpack}}}
\newcommand{\toRat}{\ensuremath{\mathrm{toRat}}}
\newcommand{\isnan}{\ensuremath{\mathrm{isNaN}}}
\newcommand{\isinf}{\ensuremath{\mathrm{isInf}}}
\newcommand{\iszero}{\ensuremath{\mathrm{isZero}}}

% Link a SymFPU commit hash (third-party repo).
\newcommand{\symfpucommit}[1]{\href{https://github.com/martin-cs/symfpu/commit/#1}{\texttt{#1}}}

% Table cell markers: a filled check for yes, a sans X for no, a parenthesized
% check for partial, each on a colored cell (green / red / yellow).
% The paired* colors are provided globally by macros.sty.
\providecolor{tablePartialYellow}{HTML}{ffe9a6}
\providecolor{markQuestionAmber}{HTML}{c47f00}
\providecolor{tableRowLight}{gray}{0.96}
\providecolor{tableRowDark}{gray}{0.88}
\providecolor{tableGroupHead}{gray}{0.80}
\newcommand{\cmark}{\cellcolor{pairedThreeLightGreen}\ding{51}}
\newcommand{\xmark}{\cellcolor{pairedFiveLightRed}{\sffamily X}}
\newcommand{\pmark}{\cellcolor{tablePartialYellow}(\ding{51})}
\newcommand{\namark}{\cellcolor{tableRowDark}\textcolor{gray}{\footnotesize NA}}
% Caption-safe swatches (\cellcolor only works inside a tabular).
\newcommand{\cmarkbox}{\colorbox{pairedThreeLightGreen}{\ding{51}}}
\newcommand{\xmarkbox}{\colorbox{pairedFiveLightRed}{\sffamily X}}
\newcommand{\pmarkbox}{\colorbox{tablePartialYellow}{(\ding{51})}}
\newcommand{\namarkbox}{\colorbox{tableRowDark}{\textcolor{gray}{\footnotesize NA}}}

% Shared pipeline-figure styles: packed values are blue, values in the
% exact/unpacked world are green. Used by the SMT-LIB recipe figure of
% \cref{sec:fp:spec} and the invariant-chain figures of \cref{sec:fp:design}.
\tikzset{
  fpbox/.style={rounded corners=2pt, align=center, minimum height=0.95cm,
                minimum width=1.9cm, inner sep=3pt, font=\footnotesize},
  fpboxwide/.style={rounded corners=2pt, align=center, minimum height=0.4cm,
                minimum width=6cm, inner sep=3pt, font=\footnotesize},
  packedbox/.style={fpbox, fill=pairedOneLightBlue!70},
  packedboxwide/.style={fpboxwide, fill=pairedOneLightBlue!70},
  unpackedbox/.style={fpbox, fill=pairedThreeLightGreen!70},
  unpackedboxwide/.style={fpboxwide, fill=pairedThreeLightGreen!70},
  parrow/.style={->, thick},
}
\newcommand{\bty}[1]{{\ttfamily\scriptsize\textcolor{black!75}{#1}}}
% The per-operation pipeline row, reused inline in \cref{sec:fp:design} next to
% each operation's description. \chainrow{op symbol}{intermediate role}{intermediate type}
\newcommand{\chainrow}[3]{%
\resizebox{\textwidth}{!}{%
\begin{tikzpicture}
\node[packedbox]   (p1) at (0,0)    {inputs $x, y$\\[2pt]\bty{PackedFloat e s}};
\node[unpackedbox] (u)  at (3.5,0)  {unpacked inputs\\[2pt]\bty{EUnpackedFloat}};
\node[unpackedbox] (w)  at (7.4,0)  {#2\\[2pt]\bty{#3}};
\node[unpackedbox] (r)  at (11.2,0) {rounded result\\[2pt]\bty{EUnpackedFloat}};
\node[packedbox]   (p2) at (14.6,0) {packed output $z$\\[2pt]\bty{PackedFloat e s}};
\draw[parrow] (p1) -- node[above]{$\unpack$} (u);
\draw[parrow] (u)  -- node[above]{$#1$} (w);
\draw[parrow] (w)  -- node[above]{$\mathrm{round}$} (r);
\draw[parrow] (r)  -- node[above]{$\pack$} (p2);
\end{tikzpicture}}%
}

% Generated, prefix-named per-family stats (produced upstream, copied in by
% chapters/floating-point/plots/sync.py).
\input{chapters/floating-point/plots/smtlib-rand/smtlib-stats}       % prefix: SmtlibRand (RQ1 cross-family SMT-LIB sample)
\input{chapters/floating-point/plots/instcombine-small/smtlib-stats} % prefix: InstcombineSmall
% Machine specs (\FpMachine*, \FpSystemSpecs*) and run config (\FpConfig*) for
% the experimental setup. Both RQ1 (smtlib-rand) and RQ2 (instcombine-small)
% were run with this same config. (The thesis copies carry an Fp prefix.)
\input{chapters/floating-point/plots/latest/triple}
\input{chapters/floating-point/plots/latest/config}
% RQ1: SMT-LIB (smtlib-rand: stratified cross-family sample, sat+unsat), ours vs Bitwuzla
\newcommand{\SmtLibNumSolvedOurs}{\SmtlibRandNumSolvedOurs}
\newcommand{\SmtLibNumSolvedBitwuzla}{\SmtlibRandNumSolvedBitwuzla}
\newcommand{\SmtLibNumTotal}{\SmtlibRandNumProblemsTotal}
\newcommand{\SmtLibNumSatOurs}{\SmtlibRandNumSatOurs}
\newcommand{\SmtLibNumUnsatOurs}{\SmtlibRandNumUnsatOurs}
\newcommand{\SmtLibNumErrorsOurs}{\SmtlibRandNumErrorsOurs}
\newcommand{\SmtLibNumTimeoutOurs}{\SmtlibRandNumTimeoutOurs}
\newcommand{\SmtLibNumCheckedOurs}{\SmtlibRandNumCheckedOurs}
\newcommand{\SmtLibNumDisagreementsOurs}{\SmtlibRandNumDisagreementsWithExpectedStatusOurs}
\newcommand{\SmtLibGeomeanTimeElapsedRatioOB}{$\SmtlibRandSpeedupBitwuzlaOverOurs\times$}
\newcommand{\KernelOverheadRatio}{$\SmtlibRandSpeedupOursNokernelOverOurs\times$}
% RQ2: small FP, ours vs exhaustive enumeration (instcombine-small)
\newcommand{\SmallFpNumSolvedOurs}{\InstcombineSmallNumUnsatOurs}
\newcommand{\SmallFpNumSolvedOursNoKernel}{\InstcombineSmallNumUnsatOursNokernel}
\newcommand{\SmallFpNumSolvedNaive}{\InstcombineSmallNumUnsatExhaustiveenumeration}
\newcommand{\SmallFpNumTotal}{\InstcombineSmallNumProblemsTotal}
\newcommand{\SmallFpNaiveTimeouts}{\InstcombineSmallNumTimeoutExhaustiveenumeration}
\newcommand{\SmallFpSpeedupOKernel}{$\InstcombineSmallSpeedupOursOverExhaustiveenumeration\times$}
\newcommand{\SmallFpSpeedupONoKernel}{$\InstcombineSmallSpeedupOursNokernelOverExhaustiveenumeration\times$}
\newcommand{\SmallFpTimeElapsedRatioOKernel}{\SmallFpSpeedupOKernel{}}
\newcommand{\SmallFpTimeElapsedRatioONoKernel}{\SmallFpSpeedupONoKernel{}}
\newcommand{\SmallFpTimeElapsedRatioON}{\SmallFpSpeedupOKernel{} with the kernel enabled (\SmallFpSpeedupONoKernel{} without)}
\newcommand{\SmallFpGeomeanTimeFpLean}{\InstcombineSmallGeomeanTimeOurs}
% Conservative lower bound on the speedup over enumeration on timed-out
% instances: \FpConfigTimeoutSec s / \SmallFpGeomeanTimeFpLean (~44x), rounded down.
\newcommand{\SmallFpTimeoutSpeedupLB}{${>}40\times$}

%%%%%%%%%%%%%%%%%%%%%%%%%%%%%%%%%%%%%%%%%%%%%%%%%%%%%%%%%%%%%%%%%%%%%%%%%%%%%%%
%% Chapter body.
%%%%%%%%%%%%%%%%%%%%%%%%%%%%%%%%%%%%%%%%%%%%%%%%%%%%%%%%%%%%%%%%%%%%%%%%%%%%%%%

\paragraph*{Attribution.} This chapter is based on \emph{Mechanizing Floating
Point Bitblasting}~\citep{bhat2027floatingpoint}, currently in submission.
\emph{My contributions:} the original unverified bitblasting solver was written
by Tunan Shi. I implemented the bitblasting circuits for multiplication,
division, and rounding, and completed the final correctness proofs for the
addition circuit. The core data structures---\texttt{PackedFloat},
\texttt{UnpackedFloat}, and \texttt{EUnpackedFloat}---were designed in
collaboration with Abdal, after studying the design of SymFPU.

\section{Overview}
\label{sec:fp:chapter-overview}
% Two paragraphs, essentially no citations.
% P1: the question this chapter answers, where it sits in the thesis argument,
%     and the result in one plain sentence.
% P2: the approximations made to the theoretically-justified algorithm, then
%     the empirical payoff.

\lettrine{F}{loating-point} verification is central to trustworthy numerical
computing. Existing approaches fall into two camps: proof assistants that
provide precise formalizations of floating-point semantics, such as
Flocq~\citep{boldo2011flocq}, but require manual proofs; and SMT solvers such
as cvc5~\citep{barbosa2022cvc5} and Bitwuzla~\citep{niemetz2023bitwuzla} that
perform efficient, push-button reductions of floating point to bitvectors via
bitblasting with a ``soft floating-point unit.'' The latter approach relies
on the remarkable progress in bitvector SMT solving---including the verified
bitblasting infrastructure of \cref{ch:background-bv}---but the floating-point
reductions themselves are subtle, complex, and had never been verified.

We bridge this gap by presenting the first verified floating-point bitblaster:
a library and associated tactic in the Lean proof assistant that
(a) mechanizes the SMT-LIB theory of floating point,
(b) extends Lean's bitvector solver with floating-point bitblasting capabilities,
(c) proves the correctness of the bitblasting reduction relative to the SMT-LIB theory, and
(d) carefully builds the underlying circuits to support small floating-point formats.
Our evaluation shows that verified bitblasting is practical:
on the SMT-COMP benchmarks our solver runs within a \SmtLibGeomeanTimeElapsedRatioOB{} geomean slowdown of Bitwuzla,
while offering far greater correctness guarantees than before.

\section{How to Bitblast a Floating-Point Multiply}
\label{sec:fp:overview}

Consider proving that floating-point multiplication is commutative:
$\forall x\, y,\ x \times y = y \times x$.
To decide this question with a SAT solver, we negate the goal and search for a counterexample:
$\exists x\, y,\ x \times y \neq y \times x$.
A floating-point number of format such at $E5M4$ is just $1 + 5 + 4$ bits,
so if we can write the left-hand side and the right-hand side as Boolean circuits
that compute the bits of $x \times y$ and of $y \times x$ from the bits of $x$ and $y$,
the negated goal becomes a gigantic propositional formula:
does some assignment to the input bits make the two circuits disagree on an output bit?
We feed this formula to a SAT solver, where
an unsatisfiable answer proves commutativity,
while a satisfying assignment is a concrete counterexample.
Everything therefore hinges on constructing a circuit that computes $x \times y$
exactly as the IEEE semantics prescribes.
This section walks through the algorithmic structure of that multiplication circuit.
Multiplication is a good case study, as it is one of the easy floating-point operations,
and yet highlights the key steps of the bitblasting pipeline.

\begin{wrapfigure}[30]{r}{0.42\textwidth}
\centering
\small
\begin{minipage}{0.40\textwidth}
\textbf{Unpack}\\[2pt]
\quad $m_x, e_x \gets \mathrm{unpack}(x)$ \quad $m_x \in [1,2)$\\
\quad $m_y, e_y \gets \mathrm{unpack}(y)$ \quad $m_y \in [1,2)$\\[6pt]
%
\textbf{Handle Special Cases}\\[2pt]
\quad $\mathit{out} = \nan \iff \mathrm{isNaN}(x) \lor \mathrm{isNaN}(y)$\\
\quad $\phantom{\mathit{out} = \nan \iff{}} \lor\ (0 \times \pm\infty)$\\
\quad $\mathit{out} = \pm\infty \iff \mathrm{isInf}(x) \lor \mathrm{isInf}(y)$\\
\quad $\mathit{out} = \pm 0 \iff \mathrm{isZero}(x) \lor \mathrm{isZero}(y)$\\[6pt]
%
\textbf{Fixed-Point Multiply} \quad $m_z \in [1, 4)$\\[2pt]
\quad $m_z \gets m_x \times m_y$\\
\quad $e_z \gets e_x + e_y$\\[6pt]
%
\textbf{Renormalize} \quad $m_z' \in [1, 2)$\\[2pt]
\quad $m_z', e_z' \gets \mathrm{renormalize}(m_z, e_z)$\\[4pt]
%
\textbf{Round}\\[2pt]
\quad $m_\mathit{out}, e_\mathit{out} \gets \mathrm{round}(m_z', e_z')$\\[6pt]
%
\textbf{Pack}\\[2pt]
\quad $\mathit{out}_\mathit{ieee} \gets \mathrm{pack}(m_\mathit{out}, e_\mathit{out})$
\end{minipage}
\caption{Algorithmic overview of floating-point multiplication:
  unpack, handle special cases, compute via fixed-point arithmetic,
  renormalize, round, and pack.}
\label{fig:fp:algo-overview}
\end{wrapfigure}

The SMT-LIB and IEEE semantics define multiplication of floating-point numbers
by representing them as extended real numbers (reals with infinity and NaN),
computing the multiplication as real numbers, and then \emph{rounding} the result back to floating-point format,
by choosing the floating-point number that is closest to the real-number result.
This description is declarative: it does not tell us how to compute the result, but rather what the result should be.

In order to implement this multiplication in a bitblasting-based solver,
we must create an arithmetic circuit that computes the multiplication and rounding such that the result matches the SMT-LIB semantics.
To be concrete, we trace a single multiplication:
$\baseten{1.5} \times \baseten{1.5}$, written in binary as $\basetwo{1.10} \times \basetwo{1.10}$.
We work in a small format \textbf{E3M2} (a three-bit exponent and a two-bit significand), in which the number $\baseten{1.5}$ is
written in scientific notation as $\basetwo{1.10} \times 2^0$.
Recall that in the floating-point format E3M2,
we are allowed to represent numbers of the form $1.m_1m_2 \times 2^e$ where $e$ is a three-bit exponent and $m_i$ are the bits of the significand.\footnote{This discussion suppresses subnormal numbers and the sign bit.}  Thus, we write the number $\baseten{1.5}$ as $\basetwo{1.10} \times 2^0$ in this format, where the significand is $\basetwo{1.10}$ and the exponent is $0$; as a packed bit-pattern (biased exponent followed by significand) it is $\bitpattern{011\,10}$.
We will now trace the steps of the multiplication algorithm by multiplying $x = y = \baseten{1.5}$,
via the multiplication bitblasting circuit (summarized in \cref{fig:fp:algo-overview}).


\noindent\textbf{Unpack.}
We first unpack each packed floating-point number (which is the dense bit-pattern $\bitpattern{011\,10}$ that represents the number)
into a representation that tags the special values (\nan{}, $\pm\infty$, $\pm 0$) explicitly and puts
every nonzero finite number into the form $2^{e} \times 1.m_1m_2$.
Our inputs are ordinary normal numbers, and each $\baseten{1.5}$ unpacks to significand
$m_x = m_y = \basetwo{1.10}$ and exponent $e_x = e_y = 0$.

\noindent\textbf{Handle Special Cases.}
We then dispatch on the special cases of \nan{}, infinity, and $\pm 0$.
For example, $\pm 0 \times \pm\infty$ is $\nan$.
Neither of our operands is special, so we take the numeric path and may assume
from here on that both inputs are nonzero finite numbers.

\noindent\textbf{Fixed-Point Multiply.}
We compute $(m_x \times 2^{e_x}) \times (m_y \times 2^{e_y}) = (m_x \times m_y) \times 2^{e_x + e_y}$.
This gives us $e_z = e_x + e_y = 0 + 0 = 0$ and $m_z = m_x \times m_y = \basetwo{1.10} \times \basetwo{1.10} = \basetwo{10.01} = \baseten{2.25}$.
Now, our resulting number is $z = m_z \times 2^{e_z} = \basetwo{10.01} \times 2^0 = \baseten{2.25}$.

\noindent\textbf{Renormalize.}
Because $z = \basetwo{10.01} \times 2^0$ does not have a significand of the normalized form $1.b_1b_2$,
we renormalize it, to write it as $z = \basetwo{1.00\bluebox{1}} \times 2^1$.
The highlighted bit $\bluebox{1}$ is an \emph{extra} bit that falls beyond the two-bit significand;
the rounding circuit uses it to choose the final result.
See that this still represents the same number: $\basetwo{1.001} \times 2^1 = \basetwo{10.01} \times 2^0 = \baseten{2.25}$.
What we did achieve was to establish the invariant that $m_z$ is normalized.
This is used by the rounding circuit to determine the closest floating-point numbers to $m_z$,
since one can ``read off'' the closest representable numbers from the normalized representation.

\noindent\textbf{Round.}
Abstractly, the rounder takes the exact normalized representation $\basetwo{1.00\bluebox{1}} \times 2^1 = \baseten{2.25}$ together with the target
precision (a two-bit significand) and returns the closest representable floating-point number.
It does so by bracketing $\baseten{2.25}$ between a lower representable value and an upper representable value.
These are the lower bound $\basetwo{1.00} \times 2^1 = \baseten{2.0}$
and the upper bound $\basetwo{1.01} \times 2^1 = \baseten{2.5}$.
See that these can be read off from the normalized representation:
the lower bound is obtained by truncating the significand to two bits, and the upper bound is obtained by adding one to the truncated significand.
In our case, the number $\baseten{2.25}$ is \emph{exactly} the midpoint of $\baseten{2.0}$ and $\baseten{2.5}$ (the extra bit $\bluebox{1}$ is a $1$ with nothing below it).
Here, the rounding mode (round to nearest-even) breaks the tie toward the even significand, returning
$\basetwo{1.00} \times 2^1 = \baseten{2.0}$.
All of these computations can be encoded as arithmetic circuits.

\noindent\textbf{Pack.}
Now that we have the rounded result $\basetwo{1.00} \times 2^1 = \baseten{2.0}$,
we pack it back into the E3M2 bit layout, giving the packed bit-pattern $\bitpattern{100\,00}$,
the final product of the multiplication.

Since each of the above phases can be expressed as an arithmetic circuit, the entire computation can be encoded into a SAT formula.
This formula is fed into a SAT solver for proving properties about this computation.
Overall, we mechanize every one of these steps of the circuit in the Lean proof assistant,
and establish semantic equivalence with the SMT-LIB specification of floating-point semantics,
which was designed to match IEEE 754.

\subsection{Why Is the Multiplication Circuit Correct?}\label{sec:fp:overview-correctness}

While we have informally argued that the multiplication circuit is correct,
we now give a more formal overview of the correctness argument.
The SMT-LIB specification computes the multiplication over exact
numbers: embed both operands into the reals extended with $\pm\infty$ and \nan{}, multiply exactly,
and round the exact product back into the format. Our circuit instead computes with bounded bitvectors.
However, it computes the exact product of the two inputs, and therefore, no information is lost.
Recall that we computed $\baseten{1.5} \times \baseten{1.5} = \baseten{2.25}$ exactly, and then rounded it to $\baseten{2.0}$.

The intermediate product of two $s = 2$ significand floats is a $2s = 4$ significand float, which is enough to represent the exact product.
Both the SMT-LIB specification and the circuit therefore round \emph{the same rational number}.
Now that we are rounding the same number, the correctness of the circuit
reduces to the equivalence of the rounding algorithm in the SMT-LIB specification and the rounding circuit.
Both the SMT-LIB specification and the rounding circuit choose between the $\mathrm{lower}(r)$ and $\mathrm{upper}(r)$ neighbors of the exact product $r$,
with the choice determined by the location of the exact product $r$ inside the interval $[\mathrm{lower}(r), \mathrm{upper}(r)]$.


Not every operation is as well-behaved.
For example, the division $1/3 = 0.3333\ldots$ is not representable in any finite width.
For such operations, we prove that we do not need to compute the exact result,
but rather a \emph{lossy summary} of it.
The summary needs enough bits to determine the two representable neighbors of the exact result,
as well as to let us know whether the truncation error is zero.
We then prove that this lossy summary is indistinguishable
from the exact result for rounding.
Now that we have an overview of the algorithm and the correctness argument,
we will explain in greater detail the prose of the SMT-LIB specification,
and how we mechanize it in Lean.

\section{A Full Mechanization of the SMT-LIB QF\_FP Theory}
\label{sec:fp:spec}

The IEEE 754 standard is written as prose, which leaves room for ambiguity.
The SMT-LIB floating-point theory~\cite{extrealsemantics}
addresses this by providing a model-theoretic semantics for floating-point arithmetic.
Operations are defined over extended reals (reals augmented with $\pm\infty$ and $\nan$)
with rounding specified as a function that maps an extended real
to the nearest representable floating-point value,
breaking ties according to the rounding mode~\cite{extrealsemantics}.
This gives a rigorous, unambiguous specification
that can serve as a ground truth for solver implementations,
and it is the specification that SMT solvers such as Z3~\cite{de2008z3},
cvc5~\cite{barbosa2022cvc5}, and Bitwuzla~\cite{niemetz2023bitwuzla} implement.

Three requirements shaped the design of our mechanization of the SMT-LIB semantics.
First, the mechanization must be \emph{faithful}: close enough to the SMT-LIB
text that transcription errors are unlikely.
Second, it must be \emph{provable against}: stated in terms that make the
circuit-correctness proofs tractable (\cref{sec:fp:operation-proofs}).
Third, it must be \emph{executable}, so that the specification itself can be run and differentially tested (\cref{sec:fp:testing}).
These requirements pull in different directions.
We reconcile these needs by making the semantics generic over its numeric domain.
This results in a single mechanization that meets all three requirements,
serving simultaneously as reference semantics and as test oracle.

The entire semantics follows a single recipe:
embed the packed IEEE floats into a numeric domain \texttt{R} that extends the reals with $\pm\infty$ and \nan{},
perform the operation \emph{exactly} in \texttt{R}, and round the exact result back into the format.
Our mechanized semantics of multiplication (left) is a direct transcription of this recipe (right):

\medskip
\noindent
\begin{minipage}[c]{0.5\linewidth}
\begin{lean4}[fontsize=\footnotesize]
--  Multiplication SMT-LIB semantics.
def mul 
    (M : RoundMethod (PackedFloat e s) R)
    (rm : RoundingMode)
    (x y : PackedFloat e s) : 
      PackedFloat e s :=
  -- the exact product, computed in R
  let z    := M.embed x * M.embed y
  -- IEEE sign rule: 
  --  XOR of the operand signs
  let sign := x.sign ^ y.sign
  -- round back into the format
  M.round rm sign z
\end{lean4}
\end{minipage}\hfill
\begin{minipage}[c]{0.49\linewidth}
\centering
\begin{tikzpicture}[font=\footnotesize]
\node[packedboxwide]   (in)  at (0,0)    {inputs $x, y$ \bty{PackedFloat e s} (\S\ref{sec:fp:ieee754})};
\node[unpackedboxwide] (dom) at (0,-1.5) {exact values in the numeric domain \bty{R} (\S\ref{sec:fp:domain})};
\node[unpackedboxwide] (z)   at (0,-2.9) {exact product \bty{z : R}};
\node[packedboxwide]   (out) at (0,-4.3) {rounded output \bty{PackedFloat e s}};
\draw[parrow] (in)  -- node[right]{\bty{M.embed}} (dom);
\draw[parrow] (dom) -- node[right]{exact $\times$ in \bty{R}} (z);
\draw[parrow] (z)   -- node[right]{\bty{M.round rm} (\S\ref{sec:fp:adjunction})} (out);
\end{tikzpicture}
\end{minipage}
\medskip

\noindent
Note that the definition has no special cases:
outcomes such as $0 \times \pm\infty = \nan$ are decided by the exact arithmetic of \texttt{R} itself.
The other arithmetic operations follow the same shape,
differing only in their zero-sign rule and the exact operation performed in \texttt{R}.

To make sense of this definition, we need the three concepts:
the packed floats (\texttt{PackedFloat e s}) that the operation consumes and produces (\cref{sec:fp:ieee754}),
the numeric domain \texttt{R} in which the exact product is computed (\cref{sec:fp:domain}),
and the round method \texttt{M} that maps the exact product back into the format (\cref{sec:fp:adjunction}).
The next three subsections define each of these in turn.

\subsection{Mechanized SMT-LIB Definition for IEEE-754 Floats}\label{sec:fp:ieee754}

The packed bit-level representation of floating-point numbers is fixed by the IEEE 754 standard.
Recall that IEEE 754 states that floating-point numbers are represented by three parts:
a sign $s$, exponent $e$, and significand $m$ (we write $m$ for the significand, after its historical name, the \textit{mantissa}),
where the sign is represented by a single bit and the exponent and significand are represented by unsigned integer bitvectors.
Our \texttt{PackedFloat} datatype
is a \emph{bit-exact mirror} of this layout:
\begin{lean4}
structure PackedFloat (E M : Nat) where
  sign : Bool     ex : BitVec E     m  : BitVec M
\end{lean4}
\noindent The floating-point specification encodes numbers in scientific notation ($\pm\, 1.m_1 m_2 \cdots m_M \times 2^{e}$):
the significand field $m$ stores the bits $m_1 m_2 \cdots m_M$, the exponent field $\mathit{ex}$ encodes the exponent $e$ (via a bias, below), and the sign bit stores $\pm$.
Writing $E$ and $M$ for the lengths of the exponent and significand bitvectors, a \emph{normal} number denotes:
\begin{equation}
\llbracket (s, \mathit{ex}, m) \rrbracket = (-1)^s\ 2^{\,\mathit{ex} - E_0} \left(1 + \frac{m}{2^M}\right),
\end{equation}
where the exponent field $\mathit{ex}$ is an unsigned $E$-bit integer and $E_0 = 2^{E-1} - 1$ is the \textit{exponent bias}, so that the represented exponent $e = \mathit{ex} - E_0$ may be negative.
This representation leaves a large gap between the smallest normal number $(1.0\cdots0 \times 2^{1 - E_0})$ and the number $0$.
IEEE closes the gap by \emph{stealing} the smallest exponent field, $\mathit{ex} = 0$: such floats are reinterpreted as \emph{subnormals} $0.m_1 m_2 \cdots m_M \times 2^{1 - E_0}$, dropping the implicit leading $1$ so that magnitudes shrink smoothly down to zero.

This now has a representation of finite numbers, but computations can give results that are too large to represent.
For example, the division $1/0$ should not be left undefined, so we symmetrically \emph{steal the highest exponent field}, $\mathit{ex} = \bitpattern{1\dots1}$, to hold an \emph{infinity}; we keep both $+\infty$ and $-\infty$, so that overflow toward large positive and large negative magnitudes remains distinguished. This lets us preserve the identity $1/(1/a) = a$ at the extremes, but only if we also distinguish $+0$ from $-0$: without signed zeros we would have $1/{+\infty} = 0 = 1/{-\infty}$, collapsing the two and breaking the identity.
So the sign bit stays meaningful at zero, with distinct $+0$ and $-0$.

Finally, once $\pm\infty$ is available we can write indeterminate expressions such as $+\infty + (-\infty)$, whose true value depends on the rates of convergence of the operands and cannot be recovered from the floats alone.
To represent such outcomes we \emph{launder} the highest exponent field a second time: infinity occupies only the slot $\mathit{ex} = \bitpattern{1 \dots 1}$, $m = \bitpattern{0 \dots 0}$, so we take a \emph{nonzero} significand with all-ones exponent to mean \nan{} (not-a-number).
The significand bits of a \nan{} then allow the programmer to store a payload for error diagnosis.
In summary, the value denoted by a \texttt{PackedFloat} is determined by its exponent field:
\begin{equation}
\llbracket (s, \mathit{ex}, m) \rrbracket =
\begin{cases}
(-1)^s\, 2^{1 - E_0}\,\frac{m}{2^M} & \mathit{ex} = 0,\ m \ne 0 \quad\text{(subnormal)},\\[4pt]
(-1)^s\, 0 & \mathit{ex} = 0,\ m = 0 \quad\text{(signed zero)},\\[4pt]
(-1)^s\, 2^{\,\mathit{ex} - E_0}\left(1 + \frac{m}{2^M}\right) & 1 \le \mathit{ex} \le 2^E - 2 \quad\text{(normal)},\\[4pt]
(-1)^s\, \infty & \mathit{ex} = 2^E - 1,\ m = 0 \quad\text{(infinity)},\\[4pt]
\nan & \mathit{ex} = 2^E - 1,\ m \ne 0 \quad\text{(NaN)}.
\end{cases}
\end{equation}

\paragraph{Equivalence Up to NaN.}
As noted above, the packed IEEE format admits many \nan{} bit patterns (differing in their payloads),
whereas the SMT-LIB model contains a \emph{single} \nan{} value for its packed floating-point representation.
To model this, we define the relation \texttt{EquivUptoNaN},
which relates two packed floats iff they are both bit-equal, or both \nan{}.
Our top-level theorems prove that the circuit and the SMT-LIB specification agree up to \nan{}.
This gives us the floats themselves, the first ingredient of the recipe.
Next, we define the numeric domain \texttt{R} into which they embed.

\subsection{Mechanized SMT-LIB Definition for Extended Reals}\label{sec:fp:domain}

The SMT-LIB semantics provides denotations of floating points over the extended reals,
i.e., reals extended with $\pm\infty$ and \nan{}.
The real numbers in Lean live in Mathlib, Lean's large library of formalized mathematics.
Unfortunately, this is a heavyweight dependency, and we wished to keep our core mechanization of the SMT-LIB semantics Mathlib-free.
To enable this, we instead choose to work over an arbitrary numeric domain \texttt{R},
constrained only by the typeclass \texttt{ExtendedNumber}
which supplies exact arithmetic together with $\pm\infty$ and \nan{}.
This one interface is instantiated at the extended rationals for our correctness proofs,
and can be instantiated with the extended reals from Mathlib for future developments, such as floating-point numerical analysis.

\begin{lean4}[fontsize=\small]
-- The numeric domain R: exact arithmetic, extended with $\pm\infty$ and NaN.
class ExtendedNumber (R : Type) extends Mul R, Add R, LE R, Zero R where
  isNaN  : R → Prop
  isZero : R → Prop
  -- ... sign tests and NaN-aware equality elided ...
\end{lean4}
\noindent For our purposes, we use the extended rationals $\mathbb{Q} \cup \{\pm\infty, \nan\}$ as the numeric domain \texttt{R}:
\begin{lean4}[fontsize=\small]
inductive ExtRat where | nan | inf (neg : Bool) | rat (num : Q)
-- Embedding of floats into extended rationals.
class RoundableEmbed (X R : Type) where embed : X → R
instance : RoundableEmbed (PackedFloat e s) ExtRat
\end{lean4}
\noindent The extended rationals suffice because floating-point numbers denote rational numbers:
their sums, differences, products, and quotients are again rational.
The only operation that can genuinely produce an irrational result is square root,
and we sidestep the reals there with a Dedekind-cut-like characterization,
proving that for a given rational $q$, we compute the floating-point number $x$
whose square is closest to $q$ among the squares of all floating-point numbers.
This lets the entire mechanization depend only on core Lean.
Next, we describe our semantics for rounding.

\subsection{Mechanized SMT-LIB Definition for Rounding}\label{sec:fp:adjunction}

The final ingredient of the recipe is the round method \texttt{M},
which mechanizes the abstract SMT-LIB rounding semantics generically over the numeric domain \texttt{R}.
The genericity is captured by the \texttt{RoundMethod} structure
which is parameterized by a numeric domain \texttt{R} and a floating-point representation \texttt{X},
and bundles the four operations needed to implement SMT-LIB rounding:

\begin{lean4}[fontsize=\small]
-- Data necessary to perform rounding.
structure RoundMethod (X R : Type) extends RoundableEmbed X R where
  lower, upper : R → X      -- greatest float $\leq$ r; least float $\geq$ r
  lowerHalf    : R → Prop   -- is r strictly below the midpoint of [lower r, upper r]?
  tieBreak     : R → Prop   -- is r exactly the midpoint?
  isEven       : X → Bool   -- is the significand's last bit zero?
\end{lean4}
\noindent We have (1) an embedding of \texttt{X} into \texttt{R}, used to compute the exact result of an operation in \texttt{R};
(2) the maps $\mathrm{lower}$ and $\mathrm{upper}$ that compute, for an exact value $(r : \texttt{R})$, the two representable floats that bracket it;
(3) the predicates \texttt{lowerHalf} and \texttt{tieBreak} that locate $r$ inside the interval $[\mathrm{lower}(r), \mathrm{upper}(r)]$:
whether $r$ lies strictly below the midpoint, and whether $r$ is exactly the midpoint;
and (4) \texttt{isEven}, which reports whether a significand is even, for tie-breaking.

Given all of this data, the rounding function \texttt{roundRNE} is defined by a
simple case analysis on the location of $r$ inside the interval
$[\mathrm{lower}(r), \mathrm{upper}(r)]$.
\begin{lean4}[fontsize=\small, escapeinside=||]
-- Round-to-nearest-even chooses between 'lower r' and 'upper r'.
def RoundMethod.roundRNE (M : RoundMethod X R) (sign : Bool) (r : R) : X :=
  if      M.isZero r            then signedZero sign  -- $\pm$0, sign supplied externally
  else if M.lowerHalf r         then M.lower r        -- strictly in the lower half
  else if |$\lnot$| M.tieBreak r        then M.upper r        -- strictly in the upper half
  else if M.isEven (M.lower r)  then M.lower r        -- exact midpoint: round to even
  else                               M.upper r        -- exact midpoint: round to even
\end{lean4}
\noindent The other four SMT-LIB rounding modes (RNA, RTP, RTN, and RTZ) are defined by analogous case analyses,
and a dispatcher \texttt{RoundMethod.round} selects the rounder for a given \texttt{RoundingMode};
the operation semantics, the circuits, and the correctness theorems that follow are all generic over this rounding mode.

We build two instantiations of \texttt{RoundMethod}:
first, the SMT-LIB semantics, which is noncomputable and relies on choice principles,
and second, a computable semantics defined by enumeration,
which we prove is equivalent to the noncomputable SMT-LIB semantics.
We use the computable semantics for our trustworthy exhaustive-enumeration baseline solver,
as well as for certain proofs that are more easily done by computation.
We use the abstract semantics for correctness proofs that follow from the SMT-LIB specification.

The abstract instantiation embeds \texttt{PackedFloat}s into the extended rationals,
and defines the lower and upper bounds by their Galois connection property.
For an exact value $(r : \texttt{R})$ in the numeric domain,
$(\mathrm{lower}(r) : \texttt{X})$ is the \emph{greatest} float whose embedding is at most $r$,
and $(\mathrm{upper}(r) : \texttt{X})$ is the \emph{least} float whose embedding is at least $r$.
The bounds themselves are then obtained from the axiom of choice,
where Hilbert's choice operator $\varepsilon$ picks an arbitrary value satisfying the
relation:
\begin{lean4}
-- The abstract bound: *some* lawful lower bound, given by choice.
noncomputable def smtLibLower : RoundMethod (PackedFloat e s) ExtRat where
  lower (r : ExtRat) : PackedFloat e s := epsilon (fun x => IsLawfulLower r x)
\end{lean4}
Separately, we prove that for every $r$ a lawful lower bound exists and is unique,
so $\mathrm{lower}$ is a well-defined function with $\mathrm{IsLawfulLower}(r, \mathrm{lower}(r))$;
the upper bound is symmetric.
This instantiation is not computable by design, as we mirror how the SMT-LIB standard defines the lower bound.
With the bounds defined, it remains to locate $r$ inside the interval $[\mathrm{lower}(r), \mathrm{upper}(r)]$,
i.e., to define the predicates \texttt{lowerHalf} and \texttt{tieBreak}.
SMT-LIB defines both by comparison with the format that carries one extra significand bit:
the floats of \texttt{PackedFloat e (s+1)} sit exactly at the floats of \texttt{PackedFloat e s} and at their midpoints,
so one extra bit of precision is exactly enough to resolve the position of $r$ relative to the midpoint.

The computable instantiation finds the same bounds by brute
force:
it enumerates the finitely many floats of the format and selects the closest representable values directly.
We prove the abstract and computable instantiations equivalent,
which both simplifies proofs and provides a test oracle for the circuit implementation (\cref{sec:fp:testing}).
One generic definition thus serves simultaneously as the abstract specification and as the executable test oracle,
which is exactly the executability we required of the mechanization at the start of this section.

\section{Verified Floating-Point Bitblasting Circuits for QF\_FP Operations}
\label{sec:fp:design}

Now that we have seen the abstract specification of the SMT-LIB semantics,
we describe the core algorithms of our solver that implement these operations as bitblastable circuits.
Note that we mechanize all operations except for conversions between floating-point and integer/bitvector formats and rem (\cref{tab:fp:supported-ops}). We continue sticking with multiplication as the running example, and use other operations to draw comparisons and highlight differences when necessary.


\begin{table}[t]
\centering
\scriptsize
\setlength{\tabcolsep}{3pt}
\ra{1.2}
% Two panels laid side by side to use the full text-width and avoid the tall,
% mostly-empty single column.
\newcommand{\OpTableHead}{%
  \textbf{SMT-LIB op} & \textbf{SymFPU} & \textbf{Impl.} & \textbf{Proven}
    & \shortstack{\textbf{FP8}\\\textbf{oracle}} \\}
\begin{minipage}[t]{0.49\textwidth}
\centering
\begin{tabular}[t]{@{}p{2.5cm} cccc@{}}
\OpTableHead
\rowcolor{tableGroupHead}\multicolumn{5}{@{}l@{}}{\textbf{Arithmetic}} \\
\rowcolor{tableRowLight}\texttt{fp.add}            & \cmark & \cmark & \cmark & \cmark \\
\rowcolor{tableRowDark} \texttt{fp.sub}            & \cmark & \cmark & \cmark$^{\dagger}$ & \cmark \\
\rowcolor{tableRowLight}\texttt{fp.mul}            & \cmark & \cmark & \cmark & \cmark \\
\rowcolor{tableRowDark} \texttt{fp.div}            & \cmark & \cmark & \cmark & \cmark \\
\rowcolor{tableRowLight}\texttt{fp.fma}            & \cmark & \cmark & \cmark & \cmark \\
\rowcolor{tableRowDark} \texttt{fp.sqrt}           & \cmark & \cmark & \pmark & \cmark \\
\rowcolor{tableRowLight}\texttt{fp.rem}            & \cmark & \cmark & \xmark & \cmark \\
\rowcolor{tableRowDark} \texttt{fp.roundToIntegral} & \cmark & \cmark & \xmark & \cmark \\
\rowcolor{tableGroupHead}\multicolumn{5}{@{}l@{}}{\textbf{Sign / select}} \\
\rowcolor{tableRowLight}\texttt{fp.neg}            & \cmark & \cmark & \cmark & \cmark \\
\rowcolor{tableRowDark} \texttt{fp.abs}            & \cmark & \cmark & \cmark & \cmark \\
\rowcolor{tableRowLight}\texttt{fp.min}            & \cmark & \cmark & \cmark & \cmark \\
\rowcolor{tableRowDark} \texttt{fp.max}            & \cmark & \cmark & \cmark & \cmark \\
\end{tabular}%
\end{minipage}%
\hfill
\begin{minipage}[t]{0.49\textwidth}
\centering
\begin{tabular}[t]{@{}p{2.5cm} cccc@{}}
\OpTableHead
\rowcolor{tableGroupHead}\multicolumn{5}{@{}l@{}}{\textbf{Comparison}} \\
\rowcolor{tableRowLight}\texttt{fp.eq}, \texttt{=} & \cmark & \cmark & \cmark & \cmark \\
\rowcolor{tableRowDark} \texttt{fp.lt}             & \cmark & \cmark & \cmark & \cmark \\
\rowcolor{tableRowLight}\texttt{fp.leq}            & \cmark & \cmark & \cmark & \cmark \\
\rowcolor{tableRowDark} \texttt{fp.gt}             & \cmark & \cmark & \cmark & \cmark \\
\rowcolor{tableRowLight}\texttt{fp.geq}            & \cmark & \cmark & \cmark & \cmark \\
\rowcolor{tableGroupHead}\multicolumn{5}{@{}l@{}}{\textbf{Classification}} \\
\rowcolor{tableRowLight}\texttt{fp.isNaN}         & \cmark & \cmark & \cmark & \cmark \\
\rowcolor{tableRowDark} \texttt{fp.isInfinite}    & \cmark & \cmark & \cmark & \cmark \\
\rowcolor{tableRowLight}\texttt{fp.isZero}        & \cmark & \cmark & \cmark & \cmark \\
\rowcolor{tableRowDark} \texttt{fp.isNormal}      & \cmark & \cmark & \cmark & \cmark \\
\rowcolor{tableRowLight}\texttt{fp.isSubnormal}   & \cmark & \cmark & \cmark & \cmark \\
\rowcolor{tableRowDark} \texttt{fp.isNegative}    & \cmark & \cmark & \cmark & \cmark \\
\rowcolor{tableRowLight}\texttt{fp.isPositive}    & \cmark & \cmark & \cmark & \cmark \\
\rowcolor{tableGroupHead}\multicolumn{5}{@{}l@{}}{\textbf{Conversion}} \\
\rowcolor{tableRowLight}\texttt{to\_fp} (from bits) & \cmark & \cmark & \xmark & \cmark \\
\rowcolor{tableRowDark} \texttt{fp.to\_ubv}        & \cmark & \xmark & \xmark & \namark \\
\rowcolor{tableRowLight}\texttt{fp.to\_sbv}        & \cmark & \xmark & \xmark & \namark \\
\end{tabular}%
\end{minipage}
\caption{Supported SMT-LIB floating-point operations.
\cmarkbox~yes; \xmarkbox~no; \pmarkbox~partial;
\namarkbox~not applicable.
\textbf{Proven.} proven correct against SMT-LIB semantics for all rounding modes.
\textbf{Impl.} implemented and tested (\cref{sec:fp:testing}).
\textbf{FP8 oracle.} operation results agree with the PyMPF oracle for FP8 (\cref{sec:fp:fptg-oracle}).}
\label{tab:fp:supported-ops}
\end{table}

The circuits reuse the standard four-stage pipeline of a hardware floating-point unit as explained previously: unpack, operate, round, and pack:

\begin{center}
\chainrow{op}{operation result}{UnpackedFloat e s}
\end{center}
\noindent This is the same structure that SymFPU uses to encode floating point into bitvectors, but what is new is that we mechanize it and prove each stage correct against the SMT-LIB semantics.
The core structure is that of the operation itself, and the rounder that is applied to its result.
First, we describe the core datatype, that of an extended unpacked floating-point number,
which is the intermediate representation that the operation and the rounder operate on.
This representation is produced from a packed float by the unpacker, and is consumed by the packer to produce a packed float.
Next, we describe the rounder first (\cref{sec:fp:rounding}), because it is the most complex stage, and it fixes the contract for the operation implementation.
To round correctly, the rounder needs either the exact result of an operation, or an inexact result together with a bit that records whether any nonzero information was discarded (the \emph{sticky bit}).
Next, we describe the precise idea behind the rounding circuit, and then the operations themselves (\cref{sec:fp:operations}),
showing how each one produces a widened intermediate that meets this contract, either exactly or by setting the sticky bit.
Overall, we mechanize and prove correct the entire floating-point pipeline,
from packed inputs to packed outputs, for all operations as indicated in \cref{tab:fp:supported-ops}.


\subsection{Extended Unpacked Floating-Point Representation for Circuits}

Recall that our circuits internally operate on (extended) unpacked floating-point numbers,
which are numbers that are canonically in scientific notation.
They represent numbers of the form $\pm 1.m_1m_2 \dots m_N \times 2^{k}$, where $k$ is the magnitude and $m_1 m_2 \dots m_N$ is the significand,
or \nan{} or $\pm\infty$, or $\pm 0$.
Crucially, the numbers are stored as two bitvectors, one that is $\langle 1 m_1 m_2 \dots m_N \rangle$ and the other that is the exponent $k$ as a signed integer,
and this acts as a canonical representation of the number, without the circuit having to worry about subnormals and other details of the packed representation.
This allows the circuits to easily perform arithmetic on the significand and exponent.
Concretely, it is a bundle of a sign bit, a signed exponent, and an integer significand, all as bitvectors:
\begin{lean4}[fontsize=\small]
structure UnpackedFloat (e s : Nat) where
  sign : Bool    ex : BitVec e    sig : BitVec s
\end{lean4}
\noindent The significand \texttt{sig} carries the hidden bit explicitly, so for a normalized number its most significant bit (bit $s-1$) is $1$, and the binary point sits immediately after it.
Treating \texttt{sig} as an integer and folding the scaling factor $2^{-(s-1)}$ into the exponent means that alignment, addition, and rounding reduce to bit shifts, integer addition, and MSB tests, with no fractional arithmetic.
Subnormal packed numbers are normalized during unpacking, which is why \texttt{e} is chosen wide enough to hold both the smallest normalized-subnormal exponent and every normal exponent.

This numeric payload cannot, on its own, represent the exceptional IEEE values \nan{} and $\pm\infty$. The extended representation therefore pairs it with an explicit classification tag:
\begin{lean4}[fontsize=\small]
inductive State | NaN | Infinity | Number
-- Extended unpacked float: a classification tag plus a numeric payload.
structure EUnpackedFloat (e s : Nat) where
  state : State    num : UnpackedFloat e s
\end{lean4}
\noindent When \texttt{state} is \texttt{Number} the payload \texttt{num} holds a valid unpacked float (with $\pm 0$ recorded as a \texttt{Number} whose significand is zero); the sign of an infinity or zero is likewise read off \texttt{num.sign}.
Separating the exceptional states from the numeric payload keeps illegal bit patterns unrepresentable and lets each operation handle the \nan{}/$\infty$/zero cases explicitly before performing uniform arithmetic on finite numbers, exactly as SymFPU does.

\emph{Unpacking.} To unpack, the packed bitvector is split into its sign, exponent, and significand fields,
and the real value of the exponent is computed by subtracting the bias.
This gives us an unpacked floating-point number that represents a value of the form $1.m_1 m_2 \dots m_N \times 2^{e}$ for all numbers.
\emph{Packing.} Not all unpacked floats can be packed back, since the unpacked floating-point range is larger than the packed floating-point range.
Therefore, we have a predicate \texttt{Packable}\footnote{We prove that \texttt{Packable} characterizes exactly the image of unpacking: a \texttt{Packable} float packs and unpacks back to itself, and conversely, the unpacking of any packed float is \texttt{Packable}.} that checks whether an unpacked float can be packed back into the packed representation.


\subsection{Implementing Rounding with Guard and Sticky Bits}\label{sec:fp:rounding}

Now that we know how to unpack and pack extended unpacked floating-point numbers, we can describe the rounding stage of the pipeline.
The rounder, as input, receives an extended unpacked float of higher precision (the source precision), and must produce a packed float of lower precision (the target precision).
To round a number accurately, we need two bits of extra precision beyond the target format: a \emph{guard} bit and a \emph{sticky} bit.
The guard bit is the first bit beyond the target precision, and the sticky bit is the logical OR of all bits beyond the guard bit.
Intuitively, the guard bit tells us whether we are above or below the midpoint of the two nearest representable numbers,
and the sticky bit tells us whether we are exactly at the midpoint or not (\cref{fig:fp:guard-sticky}).
By finding the correct rounded result, the rounder guarantees that the output is representable in the target format.
This maps overflow to $\pm\infty$ or the largest finite value.
Similarly, it maps underflow into the smallest subnormal integer, or zero, as the rounding mode directs.



\begin{figure}[htb]
\centering
\resizebox{\textwidth}{!}{%
\begin{tikzpicture}[font=\footnotesize]
%% ---- top panel: the rational line with the grid of packed floats ----
\draw[->, thick] (-0.3,0) -- (13.3,0) node[right] {$\mathbb{Q}$};
\foreach \x in {0.5,1.3,2.1,2.9,4.5,6.1,9.3,12.5}
  \fill[pairedTwoDarkBlue] (\x,0) circle (2.2pt);
\node[below=6pt, pairedTwoDarkBlue, align=center] at (1.3,0) {\scriptsize packed floats};
\node[below=6pt, gray] at (2.9,0) {\scriptsize $2^k$};
\node[below=6pt, gray] at (6.1,0) {\scriptsize $2^{k+1}$};
% the exact result and its two bounds
\fill[pairedSixDarkRed] (7.4,0) circle (2.4pt);
\node[pairedSixDarkRed, above=14pt] at (7.4,0) {exact result $r$};
\draw[pairedSixDarkRed] (7.4,0.55) -- (7.4,0.1);
\draw[->, thick, bend right=30] (7.2,0.22) to node[above left=-2pt]{$\mathrm{lower}(r)$} (6.2,0.22);
\draw[->, thick, bend left=25]  (7.6,0.22) to node[above right=-2pt]{$\mathrm{upper}(r)$} (9.2,0.22);
%% ---- zoom cone into the gap around r ----
\draw[densely dashed, gray] (6.1,-0.14) -- (0.6,-2.02);
\draw[densely dashed, gray] (9.3,-0.14) -- (12.6,-2.02);
\node[gray] at (6.6,-1.1) {\scriptsize zoom into the gap around $r$};
%% ---- bottom panel: guard and sticky locate r inside the gap ----
\fill[pairedOneLightBlue!45]    (0.6,-2.56) rectangle (6.6,-2.04);
\fill[pairedThreeLightGreen!55] (6.6,-2.56) rectangle (12.6,-2.04);
\draw[thick] (0.6,-2.3) -- (12.6,-2.3);
\fill[pairedTwoDarkBlue] (0.6,-2.3) circle (2.6pt);
\fill[pairedTwoDarkBlue] (12.6,-2.3) circle (2.6pt);
\draw[thick] (6.6,-2.62) -- (6.6,-1.98);
\node[left=4pt]  at (0.6,-2.3)  {$\mathrm{lower}(r)$};
\node[right=4pt] at (12.6,-2.3) {$\mathrm{upper}(r)$};
\node[above=2pt] at (6.6,-1.98) {\scriptsize midpoint};
% four cases for r, each with its (guard, sticky) values
\draw[pairedSixDarkRed, thick] (0.6,-2.3) circle (4.2pt);
\fill[pairedSixDarkRed] (3.6,-2.3) circle (2.4pt);
\fill[pairedSixDarkRed] (6.6,-2.3) circle (2.4pt);
\fill[pairedSixDarkRed] (9.6,-2.3) circle (2.4pt);
\node[align=center] (c0) at (1.0,-3.95)  {$r = \mathrm{lower}(r)$ \textbf{exactly}\\guard $= 0$, sticky $= 0$\\nothing to round};
\node[align=center] (c1) at (4.4,-3.95)  {$r$ \textbf{in the lower half}\\guard $= 0$, sticky $= 1$\\RNE rounds to $\mathrm{lower}(r)$};
\node[align=center] (c2) at (7.7,-3.95)  {$r$ \textbf{exactly in the middle}\\guard $= 1$, sticky $= 0$\\tie: round to even};
\node[align=center] (c3) at (11.0,-3.95) {$r$ \textbf{in the upper half}\\guard $= 1$, sticky $= 1$\\RNE rounds to $\mathrm{upper}(r)$};
\draw[->] (c0) -- (0.6,-2.68);
\draw[->] (c1) -- (3.6,-2.62);
\draw[->] (c2) -- (6.6,-2.68);
\draw[->] (c3) -- (9.6,-2.62);
\end{tikzpicture}%
}
\caption{See how the abstract SMT-LIB rounding can be implemented in terms of guard and sticky bits.
\emph{Top:} the packed floats form a finite grid among the rationals.
For an exact result $r$, $\mathrm{lower}(r)$ is the greatest float below it
and $\mathrm{upper}(r)$ the least float above it.
The rounding mode is a policy for choosing between these two candidates.
\emph{Bottom:} zooming into the gap, the guard and sticky bits are two bits
of information locating $r$: the guard bit tells us whether we are in the lower or upper half of the interval.
The sticky bit answers whether we are exactly at a boundary or slightly above it.}
\label{fig:fp:lower-upper}
\label{fig:fp:guard-sticky}
\end{figure}


Next, we describe the guard and sticky bits in more detail,
and the contract the rounder expects of the operations that use it.

\subsubsection{Inexact Results and the Sticky Compression Contract}\label{sec:fp:sticky}

While multiplication is exact, recall that division need not be exact,
since computing $1/3 = \baseten{0.33\dots} = \basetwo{0.010101\dots}$ requires an infinite number of bits.
This is a shortcoming of the binary representation, not of the circuit:
the exact quotient is a rational number, but it is not representable as a finite binary fraction.
However, the circuit can still compute a finite approximation that is guaranteed to round correctly.
For example, if we want to compute the result of $1/3$ rounded to 3 bits, we can compute the first 4 bits of the quotient: $\basetwo{0.0101}$.
Then, we need to know if this representation is exact or not. In this case,
it is not exact, and thus the sticky bit is set to $1$ to indicate that there are more bits beyond the 4th bit.
Contrast this with $5/16 = \basetwo{0.0101}$. This is exact, and thus the sticky bit is set to $0$.
Returning to $1/3$: the first 3 bits are the rounded result ($\basetwo{0.010}$).
The guard bit is $1$, which indicates that the result is at or above the midpoint of the two nearest representable numbers.
The sticky bit is also $1$, which indicates that the result is above the midpoint, and thus the result is rounded up to $\basetwo{0.011} = 3/8$.
A theorem, called \emph{sticky compression}, proves why it is safe to discard the low bits of an inexact result, and to collapse
this information into a single sticky bit.

\begin{definition}[Sticky compression]\label{def:fp:sticky-compression}
At precision $k$, the value $q$ is a sticky compression of the exact
result $r$, written $\mathrm{StickyEquiv}_k(q, r)$, iff either $q = r$,
or $q$ is an odd multiple of $2^k$ and $|r - q| < 2^k$.
\end{definition}

This is exactly the contract of a guard-and-sticky implementation:
either no low bits were discarded, or the kept value has its lowest bit forced to $1$ (the sticky OR)
and the true result lies strictly within $2^k$ of the kept value.
The master lemma
then states that rounding cannot tell the difference.
We prove a theorem that connects the bit-level rounder to the SMT-LIB rounder,
which is then reused by all operations that need rounding:

\begin{theorem}[Rounding Bridge]\label{thm:fp:rounding-bridge}
Let $x$ be a normalized unpacked float carrying at least two more
significand bits and sufficient exponent range compared to a target format $(e, s)$.
Then rounding $x$ into $(e, s)$ with the bit-level rounder agrees, up to the
choice of~\nan{}, with rounding the exact rational value of $x$
with the SMT-LIB rounder, for every SMT-LIB rounding mode.\footnote{%
\texttt{UnpackedFloat.toExtRat\_round\_Rel\_smtLibRound} in our development,
which dispatches on the rounding mode to five per-mode bridge lemmas (e.g.,~\texttt{\ldots\_of\_RNE}).}
\end{theorem}

The precondition of two extra bits of precision in the input is the guard-and-sticky requirement.
The proof connects the two rounders by connecting the predicates they use to decide which side of the midpoint the value lies on,
as well as the computation of the lower and upper representable values.
The top $s+1$ bits of $x$ determine $\mathrm{lower}$ and $\mathrm{upper}$,
and the combination of guard bit and sticky bit determines whether $x$ is below, above, or exactly at the midpoint (\cref{fig:fp:guard-sticky}).
The case analysis further covers overflow to $\pm\infty$ or the largest finite value,
as well as underflow that produces the minimum subnormal or rounds to zero.
In all cases, the bit-level rounder and the SMT-LIB rounder agree on the rounded result.

\subsection{Implementing QF\_FP Operations as Circuits with Exact Rounding}\label{sec:fp:operations}

Now that we understand the contract that rounding expects,
we describe how each operation's circuit computes a widened unpacked float.
First, special cases are handled at the \texttt{EUnpackedFloat} level before the core circuit is invoked
(e.g., adding $+\infty$ and $-\infty$ produces $\nan$).
After dispatching special cases, the core circuit computes an intermediate result that is either exact or inexact with a sticky bit.

\begin{center}
\chainrow{\times}{intermediate product}{UnpackedFloat (e+1) (2s)}
\end{center}

\noindent\textbf{Multiplication} is \textcolor{pairedFourDarkGreen}{\textbf{exact}}.
The product of two $s$-bit significands fits in exactly $2s$ bits, so the schoolbook multiplication circuit loses no information.
Since both significands lie in $[1,2)$, the product lies in $[1,4)$ and the normalizing shift is by at most one bit, so the $2s$-bit intermediate holds the true product and the rounder sees it exactly.
The exactness lemma is a calculation in rational arithmetic that unfolds the fixed-point interpretation of the product, and it shows that the circuit's value \emph{equals} the exact product, with both sides rounding literally the same rational.

\begin{center}
\chainrow{\div}{intermediate quotient}{UnpackedFloat (e+1) (s+2)}
\end{center}

\noindent\textbf{Division} is \textcolor{pairedSixDarkRed}{\textbf{inexact by necessity}}.
An exact quotient is in general not representable as a finite binary fraction at any width: just as $1/3 = \baseten{0.333\ldots}$ does not terminate in decimal, $1/10 = \basetwo{0.000110011\ldots}$ does not terminate in binary.
The divider therefore keeps $s+2$ quotient bits and derives the sticky bit from the remainder, where the discarded tail is zero iff the remainder is zero.
This meets the rounder's contract with an inexact result and a correct sticky bit.

\begin{center}
\chainrow{+}{intermediate sum}{UnpackedFloat (e+1) (s+2)}
\end{center}

\noindent\textbf{Addition} is \textcolor{pairedSixDarkRed}{\textbf{inexact by choice}}.
When the operands have exponents differing by $d$,
the exact sum needs $s + d$ significand bits,
and $d$ can be as large as the entire exponent range, which is $2^{E} - 1$ for an $E$-bit exponent.
Sizing the adder to this worst case would require hundreds of bits and is infeasible to bitblast,
and is furthermore unnecessary for correct rounding,
as the guard-and-sticky lemma shows that only two extra bits are needed.
The addition circuit therefore computes an inexact result on purpose for performance.

Finally, some operations such as sign, comparison, min/max, and classification, require no rounding at all.
Negation and absolute value only touch the sign, so their correctness proofs need no rounding.
The comparison predicates work directly at the packed level.
Now that we understand the contract of the rounder and how operations are expected to meet it,
we finally survey the core implementation of each operation and its correctness proof.

\subsubsection{The Per-Operation Correctness Theorem Structure}\label{sec:fp:operation-proofs}

Each verified operation has one top-level theorem that states that the circuit and SMT-LIB specification agree up to \nan{} equivalence.
For example, for multiplication, we state:\footnote{%
\texttt{Fp.mul\_eq\_mul} in our development.
Its analogues for addition, division, and fused multiply-add are proved
the same way.}
\begin{lean4}
theorem mul_eq_mul (a b : PackedFloat e s) (rm : RoundingMode)
    (he : 3 ≤ e) (hs : 1 ≤ s) :
  (SmtLibFunctions.mul rm a b).EquivUptoNaN (PackedFloat.mul rm a b)
\end{lean4}
\noindent The left-hand side is the SMT-LIB specification (\cref{sec:fp:spec}):
embed into the extended rationals, multiply exactly, round.
The right-hand side is the bit-level circuit (\cref{sec:fp:operations}).
The proof structure mirrors the pipeline of the floating-point implementation:

\begin{enumerate}
\item \textbf{Dispatch Special Cases.} A case analysis on the kinds
  (\nan{}, $\pm\infty$, $\pm 0$, nonzero number) of both arguments.
  The special-value cases are discharged by simplification and decision procedures;
  the number/number case is handed to layer 2.
\item \textbf{Compute (In)Exact Result With Sufficient Precision.}
  The circuit's intermediate value \emph{is} a sticky compression of the exact
  result. This is the only operation-specific obligation.
\item \textbf{Round using guard+sticky.}
   Rounding the circuit's value with the bitblasting rounder agrees with rounding the exact result with the specification's rounder,
   given that the guard and sticky bits are correct.
\end{enumerate}
\noindent The reasoning about rounding is generic over all operations, and is proved once for all formats.
The computation of the result is a per-operation calculation over unpacked floats
that is sufficiently close to the exact result that the rounder is capable of rounding correctly.

\subsection{Summary of the Verified Circuit Design}

Overall, we mechanize and prove correct most of the SMT-LIB floating-point operations
with this approach (\cref{tab:fp:supported-ops}).
We prove all the field operations and comparisons correct,
as well as all the relational operations.
The one operation whose specification is unsatisfying to us is that of square root:
we prove that the result of the square root operation is closer to the exact result than either of the two nearest representable numbers.
However, since we do not use the real numbers, we cannot say that the result \emph{is} the square root.
We also do not mechanize the round-to-integral and rem operations,
as their circuits are as involved as the arithmetic circuits.
We leave mechanizing them as future work, while still retaining trust by exhaustive testing against trusted oracles.


\section{From Verified Circuits to a Solver}
\label{sec:fp:solver}


\begin{figure}[htb]
\centering
% Cancelled fragments get a red background.
\newcommand{\fcancel}[1]{{\setlength{\fboxsep}{1pt}\colorbox{pairedFiveLightRed!70}{#1}}}
\resizebox{\textwidth}{!}{%
\begin{tikzpicture}[
  font=\footnotesize,
  io/.style   ={rounded corners=2pt, fill=pairedOneLightBlue!55, align=center, inner sep=2pt},
  term/.style ={rounded corners=2pt, fill=pairedOneLightBlue!30, align=left, font=\ttfamily\small, inner sep=4pt},
  code/.style ={rounded corners=2pt, fill=pairedOneLightBlue!18, align=left,
                font=\ttfamily\small, inner sep=5pt},
  stage/.style={font=\small\bfseries, align=center},
  flow/.style ={->, thick},
]
% ---- stage 1: inputs ----
\node[io] (smt) at (0,0.55)  {SMT-LIB};
\node[io] (lean) at (0,-0.55) {Lean goal state};
% ---- stage 2: IR of packed-float operations ----
\node[term] (ir) at (2.5,0) {%
x = add~a~b\\
y = div~x~1};
% ---- stage 3: divide-by-one cancellation (div marked red) ----
\node[term] (div1) at (4.9,0) {%
x = add~a~b\\
y = \fcancel{div~x~1}};
% ---- stage 4: expanded into eunpacked SSA ----
\node[code] (exp) at (7.9,0) {%
u1~= unpack~a\\
u2~= unpack~b\\
u3~= eadd~u1~u2\\
u4~= round~u3\\
x~~= pack~u4};
% ---- stage 5: bitblast ----
\node[term] (bb) at (10.6,0) {bv\_decide};

% terse stage titles (full detail lives in the caption)
\node[stage, above=6pt of ir]   {Input IR};
\node[stage, above=6pt of div1] {Cancel $x/1$};
\node[stage, above=6pt of exp]  {Expand};
\node[stage, above=6pt of bb]   {Bitblast};

% dataflow arrows
\draw[flow] (smt.east)  to[out=0,in=150] (ir.north west);
\draw[flow] (lean.east) to[out=0,in=210] (ir.south west);
\draw[flow] (ir)   -- (div1);
\draw[flow] (div1) -- (exp);
\draw[flow] (exp)  -- (bb);
\end{tikzpicture}%
}
\caption{Our bitblaster ingests SMT-LIB and Lean expressions and lowers them to
a bitblastable circuit in stages, illustrated on \texttt{(a + b) / 1}.
First, the division-by-one identity $x / 1 \to x$ is applied, leaving \texttt{a + b}.
Then, this is \emph{expanded} into an \unpack/operate/\pack{} circuit,
which is handed to \bvdecide{}~for bitblasting.}
\label{fig:fp:architecture}
\end{figure}


Assembling the verified circuits of \cref{sec:fp:design} into a complete,
front-end-equipped solver, and doing so while keeping the whole pipeline verified is our contribution.
The individual rewrite rules we apply are largely standard (many mirror Bitwuzla's normalizer, as we note below); what is novel is that each is proven correct in our mechanization, and that they compose with an Ackermannization pass into an end-to-end trustworthy solver.
We use the verified circuits of \cref{sec:fp:design} to build a complete floating-point solver
that can receive as input both Lean terms over \texttt{PackedFloat}
and a standalone executable that can receive SMT-LIB queries in the QF\_FP theory.
The pipeline consumes the problem statement, performs symbolic rewriting, and then produces
a bitblasting problem that is handed off to Lean's QF\_BV solver, \bvdecide{}~(\cref{fig:fp:architecture}).


\subsection{Rewrite Rules for Normalization}


We apply algebraic rewrites before bitblasting
to reduce the size of the resulting circuit.
The rules fall into three groups: constant folding, algebraic identities that mirror Bitwuzla's floating-point normalization,
and expansion of surface operations into unpacked circuits. These rewrites are proven correct in our mechanization.
The identities we do apply are those that hold unconditionally in floating point,
such as the fact that sign-only operations ($\mathrm{abs}$ and negation) commute
with the classification predicates and are idempotent under $\mathrm{abs}$.

\begin{wrapfigure}[10]{r}{0.38\textwidth}
\vspace{-2em}
\centering
\small
\begin{minipage}{0.38\textwidth}
\resizebox{\textwidth}{!}{%
\begin{tabular}[t]{ll}
\textbf{Name} & \textbf{Rewrite} \\
\rowcolor{tableRowLight}\textsc{abs-abs}    & $|\,|a|\,| \to |a|$ \\
\rowcolor{tableRowDark} \textsc{abs-neg}    & $|{-a}| \to |a|$ \\
\rowcolor{tableRowLight}\textsc{neg-neg}    & $-(-a) \to a$ \\
\rowcolor{tableRowLight}\textsc{min-self}   & $\min(a, a) \to a$ \\
\rowcolor{tableRowDark} \textsc{max-self}   & $\max(a, a) \to a$ \\
\rowcolor{tableRowLight}\textsc{lt-self}    & $a < a \to \texttt{false}$ \\
\rowcolor{tableRowDark} \textsc{le-self}    & $a \leq a \to \lnot\,\isnan(a)$ \\
\rowcolor{tableRowLight}\textsc{fold} & $\mathit{op}(\overline{c_1}, \dots) \to \overline{\mathit{op}(c_1, \dots)}$ \\
\rowcolor{tableRowLight}\textsc{expand}      & $\mathit{op}(\vec x) \to \pack\bigl(\mathit{op}_u(\unpack~\vec x)\bigr)$ \\
\end{tabular}
}
\end{minipage}
\end{wrapfigure}

While the set of rewrites may seem incomplete with respect to the properties of a ring,
such as lacking $0 \times x = 0$, such a rewrite is untrue in floating point.
If $x = +\infty$, then $0 \times \infty = \nan$.
So, our rewrite rules attempt to algebraically simplify as much as possible,
which is less than what one would naively expect from regular algebraic identities.
Finally, the \textsc{expand} rule performs the lowering the rest of the pipeline depends on,
where every surface operation is rewritten into an unpacked operation wrapped in $\unpack$ and $\pack$.

\subsection{Ackermannizing Partial Floating-Point Relations $\min$ and $\max$}\label{sec:fp:ackermannization}

SMT-LIB has some relations such as $\min$ and $\max$ that are \emph{partial}.
That is, they are left unspecified on certain inputs,
and a solver is free to pick any value in the output range for those inputs that satisfies the given constraints.
When their two arguments are zeros of opposite sign, the SMT-LIB specification
allows $\min(+0, -0)$ to evaluate to either $+0$ or $-0$.
Thus, $\min$ is fully specified everywhere except on
differing-signed-zero input.
Since the SMT-LIB semantics is model-based, the solver is free to resolve the tie either way,
as long as all other constraints are satisfied.

A standard technique for eliminating such under-specified (uninterpreted) applications
is \emph{Ackermannization}~\cite{ackermannization}, which replaces each application of the relation
with a fresh variable, together with congruence constraints ensuring
that equal arguments produce equal results.
In the general case this causes quadratic blowup,
as every pair of applications must be constrained for consistency.

In the case of $\min$, since the ambiguity only arises on differing-signed zeros, we can avoid this blowup.
We introduce two fresh, universally quantified boolean variables in the top-level goal.
These booleans are then threaded into every $\texttt{fp.min}$.
\[
  \min(x, y, b_{\min}, b'_{\min}) \triangleq
  \begin{cases}
    \text{(pick $x$ or $y$ per } b_{\min}) & \text{if } (x, y) = (+0, -0), \\
    \text{(pick $x$ or $y$ per } b'_{\min}) & \text{if } (x, y) = (-0, +0), \\
    \text{(standard minimum)} & \text{otherwise.}
  \end{cases}
\]
The solver is free to assign $b_{\min}, b'_{\min}$ to whichever values it needs when constructing a model.
We do this for every partial relation in the SMT-LIB floating-point theory,
incurring the cost of two boolean variables per relation rather than a quadratic blowup,
since the SMT-LIB semantics is only under-specified on a small number of corner cases involving NaN.

\subsection{A Lean Tactic and an SMT-LIB Solver Frontend}\label{sec:fp:frontends}

The same pipeline is packaged behind two frontends,
so interactive users and large-scale benchmarks exercise identical code paths.

\noindent\textbf{Lean tactic frontend.}
Within Lean, our solver is invoked as a tactic on goals over packed floats,
in the style of \bvdecide{}.
The tactic drives the pipeline described above as staged rewriting:
constant folding and rewriting, expansion into unpacked operations, and SAT bitblasting.
An external SAT solver (CaDiCaL) produces an UNSAT certificate,
which is checked by Lean's verified LRAT checker and, optionally,
replayed through the Lean kernel for end-to-end trust.

\noindent\textbf{SMT-LIB frontend.}
For large-scale solving, we provide a command-line SMT-LIB frontend that
parses \texttt{QF\_FP} problems, elaborates them into the same Lean goals,
and runs the identical pipeline.
Our frontend is based on the existing Leanwuzla\footnote{\url{http://github.com/hargonix/leanwuzla}} SMT-LIB parser and elaborator in Lean.
The frontend itself (parser, elaboration, and preprocessing invocations) is not verified.
We perform testing against a trusted oracle to have a high degree of confidence that the unverified frontend does not introduce errors.

\noindent\textbf{Exhaustive enumeration frontend.}
We also provide a second way to decide a closed \texttt{QF\_FP} goal that does not go through SAT.
Since a floating-point format is a finite type, a universally quantified goal over packed floats is decidable, and we can settle it by enumerating every input.
We synthesize the \texttt{Decidable} instance for this quantifier.
The instance destructures a \texttt{PackedFloat} into its sign bit, exponent bitvector, and significand bitvector, and reduces a goal $\forall\,(x : \texttt{PackedFloat}~e~s),\, P(x)$ to a quantifier over $\texttt{Bool} \times \texttt{BitVec}~e \times \texttt{BitVec}~s$, which the standard \texttt{Decidable} instances for \texttt{Bool} and \texttt{BitVec} settle by finite enumeration.
We provide the analogous instances for \texttt{EUnpackedFloat}, \texttt{UnpackedFloat}, \texttt{State}, and \texttt{RoundingMode}, so that quantifiers over all of our floating-point types are decidable.
Invoking \texttt{native\_decide} on the goal compiles this enumeration to native code, and either proves the goal or produces a counterexample.
The predicate $P$ is built from the SMT-LIB operations \texttt{fp.mul}, \texttt{fp.div}, \texttt{fp.add}, \texttt{fp.fma}, and the comparisons, which the enumerator evaluates using our computable floating-point circuits.
Each circuit is proven equal to the SMT-LIB semantics (\cref{sec:fp:operation-proofs}), so every value the enumerator computes agrees with the specification.
This is what makes exhaustive enumeration a sound oracle on the small formats where Bitwuzla's experimental support is untrusted.

\noindent\textbf{UNSAT Core Proof Replay.}
For an \texttt{unsat} result, \bvdecide{}~emits an LRAT refutation certificate that Lean's
verified LRAT checker validates, and which we optionally replay through the kernel for end-to-end trust.
This reuses the implementation of Leanwuzla.

\noindent\textbf{SAT Model Reconstruction.}
For a \texttt{sat} result, we reconstruct a packed float witness,
which we then certify as a solution by evaluating the predicate on the witness.
Previously, \bvdecide{}~reconstructed models only for bitvector and boolean variables, not structures consisting
of bitvectors and booleans (even though it can verify UNSAT proofs for them).
This is because, at least in the context of theorem proving, the UNSAT core is genuinely part
of the trusted code base, while a satisfying model is typically shown as an error to the user when a theorem could not be proven.
However, for our needs, we wish to certify the SAT model for our \texttt{PackedFloat} type.
In order to do so, we extend the way \bvdecide{} reconstructs models for bitvectors and booleans.


Recall that \bvdecide{} deals with structures by abstracting over the structure-field projections.
For example, given a packed float \texttt{pf}, it abstracts over the expressions
\texttt{pf.sign}, \texttt{pf.exp}, and \texttt{pf.m} as fresh bitvector variables.
This is entirely sound, but \bvdecide{} counts a solution as potentially spurious,
since it worries that this abstraction may be lossy. That is, it treats this as a \emph{one-way} over-approximation,
and so declines to reconstruct the abstracted boolean, exponent, significand bitvectors back into a \texttt{PackedFloat}.
However, since a structure is entirely determined by its projections, this is actually a safe over-approximation.

Therefore, we extend the model reconstruction to treat projections as safe over-approximations.
Given a satisfying SAT assignment, we read off the reconstructed sign, exponent, and significand,
reassemble the corresponding \texttt{PackedFloat}, substitute it for the witness in the
original goal, and discharge the resulting closed proposition by kernel evaluation with \texttt{decide}.
A witness that fails this check is rejected, so a reported \texttt{sat} answer is never unsound.

This leaves the trusted code base with only the unverified SMT-LIB parser and elaborator,
which we mitigate by exhaustive and randomized testing. Overall, for our solver,
\texttt{unsat} is certified by LRAT proof checking,
and \texttt{sat} by a kernel check of the satisfiability of the model.
This provides trustworthy results for both outcomes.

\section{Exhaustive and Randomized Testing for Specification Fidelity}
\label{sec:fp:testing}

Our proofs of correctness are stated relative to our translation of the SMT-LIB
floating-point specification into Lean, so it is crucial that this translation of SMT-LIB into Lean is faithful,
since a subtle mistranscription would let us soundly prove a theorem about the wrong semantics.
Toward this, we exercise our mechanized semantics against independent reference implementations of floating-point arithmetic.

On small formats we test \emph{exhaustively}, over every input (and every
pair of inputs for binary operations); on the larger IEEE formats, where exhaustive
enumeration is infeasible, we randomly sample inputs.
We check that our mechanized semantics agrees with the reference implementations
on every input, and report any discrepancies as bugs in our mechanization.

\subsection{Testing the SMT-LIB Specification via Executable Semantics}

The computable instantiation of our semantics lets us execute the specification itself, not just the circuits.
We sweep all combinations of floating-point formats with
$3 \leq E, M \leq 10$ and compare the values our specification generates
against our floating-point circuits. This gives us trust that our SMT-LIB mechanization and circuits
are consistent with each other, and that our mechanization of the SMT-LIB specification is faithful to the original SMT-LIB specification.

\subsection{Differential Testing Against Lean's Float Model}

In mid-June 2026, the Lean toolchain added a reference semantics of its
floating-point type in terms of an abstract execution model\footnote{\url{https://github.com/leanprover/lean4/pull/14079}}.
This model is validated by differential testing against Berkeley TestFloat~\cite{testfloat,softfloat} vectors:
for both binary32 and binary64, arithmetic (\texttt{add}, \texttt{sub}, \texttt{mul}, \texttt{div}, \texttt{sqrt}),
comparisons, and all conversions between floats and 32/64-bit integers are checked bit-for-bit,
with NaN results compared as a class since the model produces canonical NaNs.
It is therefore a trustworthy reference implementation of Lean's own floating-point types.
While this provides executable semantics for Lean's own floating-point types,
it provides no proof of correctness against any specification.
We would like our bitblaster to be able to reason about Lean's own floating-point types. Therefore, as a stepping stone toward this,
we first check that our mechanized SMT-LIB-based circuits agree with Lean's reference model.
Eventually, we would like to prove an isomorphism between our mechanized SMT-LIB semantics and Lean's reference model,
allowing our bitblasting tactics to be used as a drop-in library to reason about Lean's own floating-point types.

For a given floating-point format $(E, M)$, we enumerate or sample \texttt{PackedFloat} bit patterns,
run each operation on both our implementation and Lean's model, and check that the
resulting packed bits are identical (treating any two NaNs as equal).
The Lean model exposes a subset of the SMT-LIB operations,
which covers addition, subtraction, multiplication, division, square root,
negation, absolute value, and the comparisons \texttt{le}, \texttt{lt}, and \texttt{beq}.

We enumerate the two FP8 formats E5M2 and E3M4 exhaustively over all inputs, and
over all input pairs for binary operations, and find that our semantics agrees with the model bit-for-bit on every operation.
For the larger IEEE formats binary16, binary32, and binary64, where exhaustive
enumeration is infeasible, we perform randomized testing and observe the same agreement.
From these experiments, we conclude that Lean's reference model is faithful to the SMT-LIB specification,
and agrees with our mechanized semantics on all tested inputs.
We would like to eventually prove a formal isomorphism between our mechanized SMT-LIB semantics and Lean's reference model,
but for now, we only have a high degree of confidence that our semantics and Lean's reference model are equivalent.

% The FloatModel differential harness lives in FpExecutableTests/AgainstFloatModel.lean;
% the exact CI configuration is test/ci/run_float_model_tests.py. That script invokes the
% `fp-executable-test-runner` executable with a `floatModel_<fmt>` test name (plus optional
% sampling flags) once per format, runs all five configurations concurrently, and exits
% non-zero if any operation mismatches the model on any sample. The configurations are:
%   e5m2   (FP8/binary8) -- exhaustive (all 2^8 inputs, all pairs for binary ops)
%   e3m4   (FP8/binary8) -- exhaustive
%   e5m10  (binary16)    -- random, 500 samples, seed 7
%   e8m23  (binary32)    -- random, 400 samples, seed 13
%   e11m52 (binary64)    -- random, 200 samples, seed 23
% Each operation is checked bit-for-bit against Float.Model under RNE rounding.
% Excluded (no format-generic model counterpart): min/max, remainder, FMA, roundToInt.

\subsection{Differential Testing Against SymFPU}

While the Lean model is a trustworthy reference implementation,
we wanted to have a second independent check of our circuits against the de facto semantics of the solver ecosystem we aim to join.
Therefore, we differentially test our circuits against SymFPU,
the C++ library that state-of-the-art SMT solvers such as cvc5 and Bitwuzla use to implement floating-point bitblasting.

Unfortunately, SymFPU (\texttt{master} branch, commit \texttt{c3acaf6}, dated May 17, 2019) structurally
cannot process small-significand formats.
Its unpacked representation widens the exponent field until even the smallest
subnormal can be normalized, and both its unpacked-float validity check and its
rounder then align this widened exponent against the significand with an
extension operation (\texttt{matchWidth}) whose precondition demands that the
significand be at least as wide as the widened exponent.
When the precondition fails, the executable backend aborts by assertion the
moment a value is \emph{unpacked}, before any arithmetic runs.
Sweeping exponent widths $3 \le E \le 8$, we find that SymFPU's
unpack--round--pack pipeline runs exactly when $M \ge E$
(or $M \ge E - 1$ at $E = 3$).
Notably, \emph{both} machine-learning FP8 formats fail this condition:
E5M2 ($2 < 5$) and E4M3 ($3 < 4$) abort on every operation.
We therefore test in the E3M4 format, which satisfies the constraint.
Even there, the division, square-root,
and round-to-integral circuits still crashed. We traced the division failure to intermediate
bitwidths in the rounder that were one bit too short and patched SymFPU to prevent it
(widening four intermediate bitvectors), but the other two had
more involved errors that we chose not to pursue.

All of the above was against \texttt{master} commit \texttt{c3acaf6} (May 17, 2019), which is where we
began. When we spoke to the SymFPU developers,
they told us that \texttt{c3acaf6} predates a reworking of the unpacker, rounder, and divider, and pointed
us to that reworked code, which we vendored into our development at commit \texttt{502cd63} (May 13, 2026).
The \texttt{master} branch we began from has not moved since \texttt{c3acaf6} in 2019.
However, SymFPU is now under active development again on \texttt{main}, which carries this rework
together with a series of rounder and float-to-bitvector conversion fixes made through 2025 and 2026.
\footnote{Commit history on \texttt{main}: \url{https://github.com/martin-cs/symfpu/commits/main}.
Representative fixes in this burst include the FMA extended-format rework \symfpucommit{e203440} and the subnormal square-root fix \symfpucommit{b3b7977} (both March 2025);
the rounder-precondition fix \symfpucommit{d987ad4} (January 2025) and fixed-position rounder fix \symfpucommit{b213ade} (April 2026);
and the float-to-bitvector conversion fixes \symfpucommit{40bdec0} (April 2026) and \symfpucommit{506e5e7} (May 2026).}
We switched our differential oracle to this reworked version and wrote a small
enumerator that runs SymFPU's executable backend on every operation, rounding mode, and
format with $3 \le E \le 8, 1 \le M \le 8$, isolating each configuration in its own
process so that an assertion abort is attributed to a single format and operation.\footnote{\texttt{third-party/symfpu-experimental/stresstest.\{cpp,py\}} in our development.}
On this version the E3M4 square-root crash and the E4M3 abort are gone, and the
small-significand boundary moves from $M \ge E$ to $M \ge 3$.
One failure remained within the swept range: every format with $M \le 2$, including E5M2,
still aborts when a value is unpacked.
The arithmetic operations do run on E3M4, and we use it as our circuit-level cross-check.
Their results agree with our circuits, giving us confidence that our circuits are faithful
to the standard semantics of the solver ecosystem.

\subsection{End-to-End Oracle Testing of the SMT Frontend on FP8}\label{sec:fp:fptg-oracle}

The previous tests check the correctness of our mechanized semantics and the arithmetic circuits in isolation.
However, they do not check the end-to-end correctness of the solver pipeline,
where SMT-LIB parsing, preprocessing invocations, and problem creation
could have bugs, even if the core bitblaster is verified correct.
Therefore, we check our solver end-to-end against the FPTG SMT-LIB test suite.
This is the output of Schanda's floating-point test generator~\cite{symfpu},
which was designed to stress-test floating-point solvers and find bugs in them.
Each benchmark is a ground \texttt{QF\_FP} problem that binds constant inputs, applies one floating-point
operation under a given rounding mode, and asserts that the result equals a
precomputed constant. The expected \texttt{sat}/\texttt{unsat} result is stored
in the benchmark's (\texttt{:status}) field, and we check that our solver produces the expected status.

We use FP8 formats (E3M4 and E5M2), where the suite enumerates the floats to
serve as an exhaustive check and where existing solvers are known to be unsound.
We run our solver over the entire FP8 suite in \emph{both} configurations (Kernel and NoKernel),
and confirm that the results agree with the oracle on all inputs.
\cref{tab:fp:supported-ops} records this as the \textbf{Exhaustive FP8 oracle}
column.

This test has helped catch real bugs in our solver implementation.
For example, in implementing sqrt, we had failed to tag certain operations as unfoldable for the bitvector normalizer,
which caused the kernel-free path to produce eight spurious \texttt{sat} results on subnormal inputs.
Having an independent oracle to check our solver's answers against was crucial for catching this bug,
which is very easy to miss since it is a property of the solver's frontend,
which remains unverified, even though the bitblasting is fully verified.
Overall, this test gives us confidence that our end-to-end solver soundly implements the SMT-LIB floating-point theory.

\section{Case Study: Gaps Discovered by Mechanization}
\label{sec:fp:surfaced}

These are the gaps in the literature and implementations we discovered while mechanizing.
We discovered some software bugs in SymFPU, and one typo in the SMT-LIB specification that we have corrected in our mechanization.

\noindent\textbf{Latent Circuit Bugs.}
SymFPU's divider carried a soundness bug that was found and fixed only
recently, a bug we encountered independently while testing our own divider
against it. The fix, which carries the result exponent through a subtract-with-borrow
and weakens the divider's postcondition to a range check rather than re-extending a
fixed-width exponent, is present in the reworked \texttt{502cd63}, and our enumerator
finds no E3M4 division abort there.
Our theorem \texttt{div\_eq\_div}
establishes that the divider circuit agrees with the SMT-LIB specification for every
format satisfying $s \geq 1, e \geq 3$.

\noindent\textbf{Exact Width Bounds Necessary for Correctness.}
SymFPU's unpacker and rounder implicitly assume that the significand is at
least as wide as its internally-widened exponent (in practice, $M \geq E$ on
\texttt{c3acaf6}, and $M \geq 3$ on the reworked \texttt{502cd63}), and abort by
assertion on formats where the assumption fails (\cref{sec:fp:testing}).
Our per-operation theorems (\cref{sec:fp:operation-proofs}) instead carry two side
conditions, $s \geq 1$ and $e \geq 3$, with no relation between $s$ and $e$: the same
hypotheses ($\mathtt{0 < s}$, $\mathtt{1 < e}$, $\mathtt{2 < e}$) appear in
\texttt{add\_eq\_add}, \texttt{sub\_eq\_sub}, \texttt{mul\_eq\_mul}, and
\texttt{div\_eq\_div}. Both machine-learning FP8 formats, E5M2 and E4M3, satisfy these
conditions, so our proofs cover them.

\noindent\textbf{Typo in the SMT-LIB specification.}
While transcribing the SMT-LIB specification, we discovered a typo in its
Table~6: the last invocation of the rounding helper uses $tb$ where it should
use $\overline{v}$.
The error was caught mechanically, when Lean produced a type error on the verbatim
transcription, since $tb$ and $\overline{v}$ have different types in our
formalization.
Type-checking provides a consistency check on prose
specifications, catching errors that are easy to overlook in manual review.



\section{Quantitative Evaluation of Solver Performance}
\label{sec:fp:evaluation}

Two questions drive the evaluation. The first is how our solver compares against
Bitwuzla, the state-of-the-art floating-point solver---what, in other words, the
overhead of a mechanized implementation is against an unverified one. The second
concerns the small floating-point formats, whose soundness in current solvers is
unknown and which matter because of their use in machine-learning workloads: how
much faster is our solver than exhaustive enumeration there? In theory, since
small formats such as FP8 have only $2^8$ values, enumeration is a feasible, if
not fast, strategy.

\subsection{Datasets}

\noindent\textbf{SMT-COMP.}
We evaluate on the \texttt{QF\_FP} (quantifier-free floating-point) division of the
2025 release of the SMT-LIB non-incremental benchmarks~\cite{smtlib2025benchmarks},
the dataset used by the annual solver competition SMT-COMP.
The division is dominated by a single family: the \emph{wintersteiger} suite of
machine-generated verification obligations for the IEEE 754 operations accounts for
$39{,}994$ of the division's $40{,}406$ files ($\approx\!99\%$),
so a uniformly random sample would be almost entirely wintersteiger and would
overlook the smaller, hand-written, real-world families.
We therefore perform stratified random sampling of the division.
We take $K=4000$ problems from each family. This results in random sampling from the dominant wintersteiger family,
and takes all the problems from all other families, which are small enough to fit within the cap.
Overall, this gives us \SmtLibNumTotal{} problems.
This exercises both the \texttt{sat} and \texttt{unsat} answers our solver can certify,
and serves as a coverage survey over the whole division.

\noindent\textbf{LLVM Small Floating-Point Rewrites.}
We introduce a new dataset of \SmallFpNumTotal{} floating-point rewrites from LLVM.
We construct it from the test cases of LLVM's InstCombine peephole optimizer: we take the files
containing floating-point instructions, run LLVM's O2 pass pipeline on them to obtain before/after LLVM files,
and translate the before and after into SMT-LIB, with a proof obligation that the output program refines the input program.
In generating the dataset, we keep only constant-free problems, so that reducing the input float and double formats to smaller formats
is unambiguous: if a problem contained a constant such as $2^{20}$,
it is not clear what the equivalent small-format value would be, since $2^{20}$ is not representable in E4M3 or E5M2.
This gives us a dataset of problems that are representative of the kinds of
rewrites that compilers perform, and that are also small enough to be
exhaustively enumerated in the small floating-point formats.

\subsection{Solvers}
\noindent\textbf{Bitwuzla.}
We use Bitwuzla, a state-of-the-art floating-point and bitvector solver,
as our roofline.
Stock Bitwuzla only supports the standard IEEE formats
(\texttt{Float16}, \texttt{Float32}, \texttt{Float64}, and \texttt{Float128}),
and rejects the small floating-point formats (e.g.,~E4M3, E5M2) that make up much of our benchmark suite.
We therefore build Bitwuzla from source with experimental floating-point support enabled
(\texttt{./configure.py -{}-fpexp}), which allows it to accept arbitrary formats.
Bitwuzla documents these experimental formats as ``use at your own risk'' owing to known issues in SymFPU~\cite{symfpu};
this is precisely the unsoundness on small floating-point formats that motivates both our verified solver and the exhaustive-enumeration baseline.
Bitwuzla is unverified, so the entire solver is part of the trusted code base.

\noindent\textbf{Ours (Kernel + NoKernel).}
We run our solver in two modes: one where the kernel is enabled,
which increases trust but requires the kernel to check the certificate for every answer
(replaying an LRAT refutation for \texttt{unsat} and validating a reconstructed model for
\texttt{sat}, \cref{sec:fp:frontends}),
and one where the kernel is disabled,
which is less trustworthy but measures the effectiveness of our algorithmic approach
without adding the overhead of verification.
In both modes, the trusted code base contains the frontend that parses SMT-LIB problems into our internal formulas (\cref{sec:fp:frontends}),
and, since \bvdecide{} uses \texttt{native\_decide},
both modes assume the \texttt{native\_decide} axiom, i.e., the correctness of the Lean compiler.
The NoKernel mode weakens the trust guarantee further,
since we do not replay the proof through the Lean kernel.

\noindent\textbf{Exhaustive.}
The exhaustive solver decides each query by enumerating over all possible inputs,
via the \texttt{Decidable}-instance enumeration described in \cref{sec:fp:frontends}.
This is slow, but establishes an unconditionally \emph{sound} baseline for the small, non-standard
formats where Bitwuzla's experimental floating-point support is documented as potentially unsound,
since its ground evaluation uses our computable semantics, proven correct against the SMT-LIB specification.
Its trusted code base is the same parsing frontend as our bitblasting solver,
together with the \texttt{native\_decide} axiom, since the enumeration itself runs via \texttt{native\_decide}.

\subsection{Experimental Configuration}\label{sec:fp:experimental-configuration}

\noindent\textbf{Machine.}
All experiments were run on a single machine equipped with an
\FpSystemSpecsProcessorName{} (\FpSystemSpecsCores{} logical cores
at \FpSystemSpecsClockMhz{}~MHz) and \FpSystemSpecsMemoryGb{}~GB of RAM.

\noindent\textbf{Resource limits and parallelism.}
Each solver invocation is given a wall-clock timeout of
\FpConfigTimeoutSec{}~s and a memory limit of \FpConfigMemoutMb{}~MB.
An invocation that exceeds either limit
is recorded as a timeout. To amortize the cost of the full benchmark sweep we
run \FpConfigNproc{} solver invocations in parallel.
% Each problem-solver pair is executed \FpConfigRuns{} times and we report the geomean aggregate.


\subsection{Our Solver Scales to SMT-COMP}

\begin{figure}[htbp]
\centering
\includegraphics[width=\textwidth]{chapters/floating-point/plots/smtlib-rand/smtlib-cactus.pdf}
\caption{\textbf{SMT-COMP.} Our solver stays within a \SmtLibGeomeanTimeElapsedRatioOB{} geomean slowdown of
  Bitwuzla on SMT-COMP, our stratified cross-family sample of the SMT-LIB 2 QF\_FP benchmarks, in exchange for end-to-end correctness.}
\label{fig:fp:smtlib-cactus}
\end{figure}

\cref{fig:fp:smtlib-cactus}
presents results on SMT-COMP.
With the kernel enabled, our solver certifies \SmtLibNumSolvedOurs{} out of \SmtLibNumTotal{} benchmarks,
compared to \SmtLibNumSolvedBitwuzla{} for Bitwuzla,
with a geomean slowdown \SmtLibGeomeanTimeElapsedRatioOB{}.
This dataset exercises both UNSAT and SAT instances, where \SmtLibNumUnsatOurs{} are \texttt{unsat}
(discharged by a kernel-checked LRAT certificate) and \SmtLibNumSatOurs{} are \texttt{sat}
(discharged by reconstructing and kernel-validating a floating-point model, \cref{sec:fp:frontends}).
The remaining instances split into \SmtLibNumTimeoutOurs{} timeouts and \SmtLibNumErrorsOurs{} errors.
The timeouts are dominated by the \texttt{fp.rem} family, whose circuit we bitblast but whose kernel
replay is expensive (discussed below).

The errors are from SMT-LIB input constructs outside our supported fragment,
mostly literal bitvector-to-float conversions (\texttt{(\_~to\_fp~$\dots$)}),
and some goals that fall outside the bitvector fragment after preprocessing.
In both cases our frontend reports an error rather than risking an unsound answer.
Whenever the benchmark carries a known \texttt{:status}, our solver never disagrees with it
(\SmtLibNumDisagreementsOurs{} disagreements), which lends greater credibility to the end-to-end pipeline.

Our bitblaster follows the same overall algorithmic approach as Bitwuzla,
which involves encoding floating-point operations into bitvector circuits that are then reduced to SAT.
However, our bitblaster lacks some of the clever circuits that SymFPU builds, as these special-case circuits require new proofs of correctness.
In particular, certain circuits exploit properties of floating-point arithmetic
(e.g.~subtraction of floating-point numbers can be exact, under the conditions of the Sterbenz lemma~\cite{sterbenz1974floating}),
which allows them to skip the rounding step and directly compute the final result.
This accounts for part of the slower solve times.
A second source of overhead is the bitvector backend itself:
\bvdecide{} performs verified bitblasting and emits checkable LRAT
certificates, and is itself within roughly an order of magnitude of
Bitwuzla's unverified bitvector engine on SMT-LIB~\cite{leanbv}, and our floating-point solver inherits this baseline cost.

Additionally, enabling the Lean kernel for proof replay introduces
a further overhead of approximately \KernelOverheadRatio{},
since the kernel must check the LRAT certificate produced by CaDiCaL.
Kernel checking times out on verifying the problems in the \texttt{fp.rem} family,
where we exhaust the \FpConfigTimeoutSec{}~s budget, compared to far fewer timeouts on these problems for the kernel-free configuration.
The \KernelOverheadRatio{} figure is therefore a lower bound on the true cost of kernel checking,
since each configuration's geomean is taken over the instances that it solves, so the kernel's geomean is computed only over instances that do not time out,
which excludes the \texttt{fp.rem} family that is extremely expensive to check,
since the circuits we build are not optimized.

Despite these overheads, our solver certifies the large majority of SMT-COMP
within reasonable time, demonstrating that a verified bitblaster
can scale to standard floating-point benchmarks.
On the first question, then: our solver scales to standard benchmarks within a
\SmtLibGeomeanTimeElapsedRatioOB{} geomean slowdown of the state-of-the-art Bitwuzla.

\subsection{Our Solver Is Faster Than Exhaustive Enumeration for Small Floating-Point Formats}

\cref{fig:fp:small-fp-cactus} shows the cactus plot comparing our solver
against exhaustive enumeration on the small floating-point benchmarks:
the \SmallFpNumTotal{} SMT-LIB problems lifted from LLVM InstCombine rewrites
in the E4M3, E5M2, BF16, and other small formats whose total width is at most 16 bits.
Our solver certifies \SmallFpNumSolvedOurs{} out of \SmallFpNumTotal{} benchmarks,
compared to \SmallFpNumSolvedNaive{} for exhaustive enumeration,
which additionally times out on \SmallFpNaiveTimeouts{} instances.
On the instances that both approaches complete, our bitblasting attains a
geomean speedup over exhaustive enumeration of \SmallFpTimeElapsedRatioON{}.
This severely understates the actual speedup,
as it is a geomean over only the instances enumeration finishes, and so ignores the \SmallFpNaiveTimeouts{} problems on which enumeration exhausts its
time budget, which are the hardest cases on which bitblasting wins most.
A back-of-the-envelope calculation shows that
the bitblaster solver discharges each of them in about \SmallFpGeomeanTimeFpLean{}, while enumeration
spends at least its full \FpConfigTimeoutSec{}~s budget.
The true speedup on these instances is therefore at least \FpConfigTimeoutSec{}~s
divided by \SmallFpGeomeanTimeFpLean{}, which we conservatively round down and report as \SmallFpTimeoutSpeedupLB{}.
Regardless, exhaustive enumeration is an unconditionally sound baseline on these small,
non-standard formats, and our bitblasting approach matches that soundness while improving on its performance.

We also report Bitwuzla, built with its experimental small-format support enabled (\texttt{-{}-fpexp}).
Bitwuzla solves all \InstcombineSmallNumSolvedBitwuzla{} instances with a geomean time of
\InstcombineSmallGeomeanTimeBitwuzla{}, which makes it
$\InstcombineSmallSpeedupBitwuzlaOverOurs\times$ faster than our kernel-enabled solver
(and $\InstcombineSmallSpeedupBitwuzlaOverOursNokernel\times$ faster than the kernel-free configuration),
and on this suite it never disagrees with the expected status.
Empirically, Bitwuzla performs flawlessly on these benchmarks.
However, it lacks is the guarantee, since its small-format support is documented as experimental and ``use at your own risk''
and it produces no certificate that could be checked independently.
Our claim on small formats is therefore that we are the fastest solver \emph{with a soundness guarantee}.
On the second question, our solver is faster than exhaustive enumeration by a geomean factor of \SmallFpTimeElapsedRatioON{} and by \SmallFpTimeoutSpeedupLB{} on the hardest instances,
while remaining sound on small formats where Bitwuzla is not guaranteed to be sound.

\begin{figure}[htbp]
\centering
\includegraphics[width=\textwidth]{chapters/floating-point/plots/instcombine-small/smtlib-cactus.pdf}
\caption{\textbf{Small formats.} On the small FP formats (E4M3/E5M2) from the LLVM InstCombine rewrites, our verified
  solver is faster and more robust than exhaustive enumeration (\SmallFpSpeedupOKernel{} geomean with the kernel,
  \SmallFpSpeedupONoKernel{} without), certifying as many or more instances and avoiding the timeouts that stall enumeration:
  on the \SmallFpNaiveTimeouts{} instances where enumeration exhausts its \FpConfigTimeoutSec{}~s budget, our solver
  finishes in about \SmallFpGeomeanTimeFpLean{} each, a speedup of \SmallFpTimeoutSpeedupLB{}.
  The unverified Bitwuzla (\texttt{-{}-fpexp}) solves all instances
  $\InstcombineSmallSpeedupBitwuzlaOverOurs\times$ faster than our verified solver,
  but offers no soundness guarantee on these formats.}
\label{fig:fp:small-fp-cactus}
\end{figure}

Since existing solvers like Bitwuzla are fast but not known to be sound on these non-standard formats,
our bitblaster and our exhaustive enumerator fill a previously unoccupied niche:
they are the only solvers for these small floating-point formats that come with a soundness guarantee,
and of the two, bitblasting is the faster.
Our solver outperforms exhaustive enumeration by
\SmallFpTimeElapsedRatioON{} in geomean time and by \SmallFpTimeoutSpeedupLB{} on the instances
enumeration cannot finish, providing verified results on formats where no proven-sound solver exists.

\section{Related Work: Mechanized Floating-Point Arithmetic}
\label{sec:fp:related}

The introduction (\cref{ch:intro}) placed this thesis among the approaches that bridge
automated speed and kernel-checked trust; here we compare against the technical
antecedents specific to floating-point mechanization and solving.

Floating point has undergone intense study in the formal methods community,
both in terms of mechanized semantics and verified solvers.
Therefore, we do not aim to provide a comprehensive survey of the literature
(for which we guide the reader to the floating-point handbook \cite{fphandbook}).
Instead, we highlight the most relevant prior work and contrast it with our contributions.

\paragraph*{Mechanized Floating-Point Semantics}
One of the earliest mechanized floating-point formalizations is due to \citet{harrison1999machine},
who formalized IEEE 754 arithmetic in HOL Light,
covering rounding, basic operations, and key properties such as error bounds.
This work established the paradigm of proving floating-point properties in a theorem prover,
but required substantial manual proof effort.
Harrison's formalization is aimed at \emph{interactively} proving mathematical properties of floating-point arithmetic (such as rounding-error bounds) by hand, and is not connected to an automated decision procedure.
Our goal is complementary and distinct:
rather than proving properties of floating point interactively, we mechanize the reduction from the SMT-LIB semantics of IEEE 754 down to the bit-level circuits that a SAT-based solver executes, so that a \emph{push-button}, certificate-producing solver becomes trustworthy end to end.
In short, Harrison verifies floating-point \emph{theorems} in a proof assistant, whereas we verify the floating-point \emph{solver} that discharges them automatically.

In Rocq there exists PFF~\cite{pff}, a generic library for floating-point reasoning
parameterized over an arbitrary radix and precision. Flocq~\cite{boldo2011flocq} is a spiritual successor that
develops a more comprehensive library that
unifies several earlier formalizations and provides a multi-radix, multi-precision
framework for proving floating-point algorithms.
Flocq models floating-point numbers as a subset of the reals satisfying format constraints,
rather than as packed bit-level representations.
This real-number-centric view is well-suited for numerical analysis.
Building on Flocq, the VCFloat~\cite{appel2022vcfloat} library
develops theory for bounding floating-point rounding errors.
This was later extended to VCFloat2~\cite{kellison2024vcfloat2},
which performs verified floating-point interval analysis.
Both tools provide high-level numerical error reasoning and
inherit Flocq's real-number semantics.
These efforts, like Harrison's, are aimed at interactively proving properties of floating-point arithmetic,
rather than mechanizing the SMT-LIB semantics and connecting it to an automated solver.

In Isabelle/HOL, \citet{isabellefp} provides a formal model of IEEE floating-point arithmetic
in the Archive of Formal Proofs. This formalization covers the standard formats
and operations but has not been connected to an automated solver.

In Lean, FLoPS~\cite{flops} formalizes the semantics of the upcoming
IEEE-P3109 standard for low-precision floating-point arithmetic and verifies
properties of numerical algorithms over it.
FLoPS targets the P3109 parametric format family and interactive analysis,
whereas we mechanize the SMT-LIB semantics of IEEE 754 and connect it to an
automated, certificate-producing solver.

\paragraph*{Floating-Point Bitblasting and Solvers}
The SMT-LIB semantics that we use, due to \citet{extrealsemantics}, forms the basis of the SMT-LIB floating-point theory.
Their formalization defines operations over extended reals with explicit rounding,
providing a precise, automatable semantics.
SymFPU~\cite{symfpu} is the de facto standard library for reducing floating-point
to bitvector operations, and is used internally by both cvc5~\cite{barbosa2022cvc5} and
Bitwuzla~\cite{niemetz2023bitwuzla}. SymFPU encodes each floating-point operation as a
bitvector circuit, but the encodings are unverified and, as we discovered,
contain implicit assumptions relating the significand and exponent widths
that exclude the small formats used in machine learning (\cref{sec:fp:testing}).
Our work mechanizes and generalizes SymFPU's algorithms, removing some of these restrictions,
and provides a machine-checked proof of correctness for the bitblasting encodings.

\citet{symfpu} describes the design of SymFPU's bitblasting encodings,
with a focus on the rounding circuit and its correctness argument.
The correctness argument is informal and relies on exhaustive testing;
we replace it with a machine-checked proof.
Our solver builds directly on \bvdecide{}, Lean's verified bitvector
bitblaster~\cite{leanbv}, which cites floating-point reasoning as a motivating
application of bitvector solving. The present work supplies that
floating-point layer, reducing the SMT-LIB floating-point theory to the
bitvector goals that \bvdecide{} discharges.
\citet{brain2019invertibility} derives invertibility conditions for floating-point formulae,
enabling quantifier elimination for floating-point constraints in SMT solvers.
\citet{cbmcfloat} presents mixed abstractions for floating-point arithmetic in the context
of CBMC~\cite{cbmc}, combining bit-precise and abstract reasoning to verify C programs with floating-point arithmetic.
Overall, we owe a great deal to the line of research that has developed the SMT-LIB floating-point theory and bitblasting encodings,
and our work mechanizes and verifies these contributions by standing on their shoulders.

\section{Discussion}
\label{sec:fp:limitations}

We have three operations that are unproven:
round-to-integral, remainder, and the conversions to \texttt{ubv} and \texttt{sbv}.
Their circuits are implemented and the solver bitblasts them, but we're missing the proof of correctness.
In a similar vein, we lack some of the sophisticated circuits that SymFPU implements,
such as the multiple fast paths for addition that exploit various subnormal/normal cases for faster symbolic bitblasting.

Closing both gaps is limited by proof effort rather than by floating-point insight.
The obligations that arise are quite nonlinear, mixing multiplication of variables, integer division and modular arithmetic, powers of two with variable exponents, and bitvector operations.
Existing automation cannot close these, so the proofs are dominated by arithmetic bookkeeping that is conceptually trivial but mechanically exhausting.
Effective proof automation for nonlinear reasoning, ideally operating uniformly across \texttt{Nat}, \texttt{Int}, \texttt{Rat}, and \texttt{BitVec},
is therefore the key enabling technology that we'd need for this at scale.
Overall, though, it seems that proving the correctness of floating point bitblasting is tedious, error-prone, and hard, but not impossibly so!
